% Resultados do eval do checkpoint Graphix-T5 (t5-base + RGAT, treinado no Spider)
% aplicado ao ScienceBenchmark (cordis, oncomx, sdss), zero-shot.
%
% Pode ser incluído num capítulo/seção do TCC via % Resultados do eval do checkpoint Graphix-T5 (t5-base + RGAT, treinado no Spider)
% aplicado ao ScienceBenchmark (cordis, oncomx, sdss), zero-shot.
%
% Pode ser incluído num capítulo/seção do TCC via % Resultados do eval do checkpoint Graphix-T5 (t5-base + RGAT, treinado no Spider)
% aplicado ao ScienceBenchmark (cordis, oncomx, sdss), zero-shot.
%
% Pode ser incluído num capítulo/seção do TCC via % Resultados do eval do checkpoint Graphix-T5 (t5-base + RGAT, treinado no Spider)
% aplicado ao ScienceBenchmark (cordis, oncomx, sdss), zero-shot.
%
% Pode ser incluído num capítulo/seção do TCC via \input{sciencebenchmark-eval-results}
% ou compilado standalone.

\section{Avaliação no ScienceBenchmark (zero-shot)}
\label{sec:sciencebenchmark-eval}

O checkpoint treinado no Spider (t5-base + RGAT, 5 épocas, 8577 exemplos de
treino, resultado no dev set do Spider: 41,9\% \emph{exact match} / 44,8\%
\emph{execution accuracy}) foi avaliado, sem nenhum re-treinamento, sobre os
299 exemplos do split de desenvolvimento do ScienceBenchmark
\cite{sciencebenchmark2023}, distribuídos em três bancos de domínios distintos:
\texttt{cordis} (projetos de pesquisa europeus), \texttt{oncomx} (biomedicina)
e \texttt{sdss} (astronomia).

\subsection{Resultado agregado}

\begin{table}[h]
\centering
\caption{Resultado do eval zero-shot do checkpoint Spider-treinado no ScienceBenchmark, por banco.}
\label{tab:sciencebenchmark-zeroshot}
\begin{tabular}{lrrrl}
\toprule
\textbf{Banco} & \textbf{n} & \textbf{exact match} & \textbf{exec} & \textbf{Exemplos avaliados} \\
\midrule
cordis            & 100 & 12,0\% & 15,0\% & 100/100 \\
oncomx            & 99  & 14,1\% & 11,1\% & 99/99 (exec: 97/99) \\
sdss              & 94  & 0,0\%  & 0,0\%  & 90/94\textsuperscript{a} \\
\midrule
\textbf{Total (agregado)} & \textbf{293} & \textbf{9,0\%} & \textbf{8,9\%} & exact: 289/293, exec: 287/293 \\
\bottomrule
\end{tabular}

\vspace{0.3em}
{\footnotesize\textsuperscript{a}4 exemplos do sdss descartados da métrica: parser
SQL do Spider (\texttt{third\_party/spider/process\_sql.py}) não suporta
parênteses agrupando condições no \texttt{WHERE}, um padrão ausente no SQL
gold do Spider mas presente no do ScienceBenchmark.}
\end{table}

Para efeito de comparação, o mesmo checkpoint avaliado no dev set do Spider
(1034 exemplos, mesmo domínio de treino) obteve 41,9\% exact match e 44,8\%
exec -- uma queda de $\sim$33 pontos percentuais ao mudar de domínio sem
re-treinamento, consistente com a expectativa de que arquiteturas
graph-aware treinadas em um conjunto de bancos pequenos e homogêneos (Spider,
$\sim$5,3 tabelas/banco em média) generalizam de forma limitada para bancos
maiores e de domínio mais específico (ScienceBenchmark, 19--25 tabelas no
\texttt{cordis}, colunas de medição numérica pouco usuais no \texttt{sdss}).

\subsection{Quebra por dificuldade}

Usando a classificação de dificuldade do avaliador oficial do Spider
(\emph{easy}/\emph{medium}/\emph{hard}/\emph{extra}, baseada em número de
componentes SQL, condições e agregações), o exact match cai a zero em
\emph{hard} e \emph{extra} nos três bancos, sem exceção:

\begin{table}[h]
\centering
\caption{Exact match por nível de dificuldade e banco.}
\label{tab:sciencebenchmark-hardness}
\begin{tabular}{lrrrr}
\toprule
\textbf{Banco} & \textbf{easy} & \textbf{medium} & \textbf{hard} & \textbf{extra} \\
\midrule
cordis (n=100) & 32,0\% (n=25) & 11,4\% (n=35) & 0\% (n=19) & 0\% (n=21) \\
oncomx (n=99)  & 42,9\% (n=21) & 13,9\% (n=36) & 0\% (n=22) & 0\% (n=20) \\
sdss (n=90)    & 0\% (n=6)     & 0\% (n=28)    & 0\% (n=20) & 0\% (n=36) \\
\bottomrule
\end{tabular}
\end{table}

O padrão observado sugere duas causas de queda distintas e sobrepostas: (i)
uma limitação de \emph{capacidade}, comum aos três bancos -- o modelo t5-base
resolve consultas simples mas não lida com consultas complexas (múltiplos
JOINs, subconsultas, agregações compostas) fora do domínio de treino; e (ii)
uma limitação de \emph{transferência de domínio}, específica do
\texttt{sdss} -- o resultado é zero mesmo nas consultas \emph{easy}, o que não
ocorre em \texttt{cordis}/\texttt{oncomx}, indicando que o vocabulário e a
estrutura de esquema do \texttt{sdss} (medições fotométricas, nomenclatura
astronômica) são suficientemente distantes do domínio de treino a ponto de
comprometer até os casos mais simples.

Uma observação metodológica reforça essa segunda leitura: o tamanho do
esquema (contagem de nós do grafo RGAT) e o banco de dados são quase
colineares nesta amostra -- \texttt{cordis}/\texttt{oncomx} sempre excedem
512 nós, \texttt{sdss} está sempre abaixo. Como \texttt{sdss} é o banco com o
\emph{menor} esquema em número de nós e ainda assim o único com resultado
nulo em todos os níveis de dificuldade, o fator determinante da queda não
parece ser o tamanho do esquema, e sim a distância de domínio/vocabulário em
relação ao Spider.

\subsection{Análise qualitativa dos erros}

Das 263 previsões incorretas (excluindo os 4 exemplos de \texttt{sdss} sem
gold parseável), 179 (68\%) foram inicialmente classificadas como ``SQL
sintaticamente inválido'' pelo parser oficial do Spider
(\texttt{third\_party/spider/process\_sql.py}) -- ou seja, o modelo nem
produziu uma consulta que o parser conseguisse interpretar como AST SQL.

\begin{table}[h]
\centering
\caption{Distribuição das categorias de erro (previsões incorretas, exact match = 0).}
\label{tab:error-categories}
\begin{tabular}{lrrrr}
\toprule
\textbf{Categoria} & \textbf{Total} & \textbf{cordis} & \textbf{oncomx} & \textbf{sdss} \\
\midrule
SQL sintaticamente inválido & 179 & 56 & 58 & 65 \\
Tabela errada                & 50  & 13 & 12 & 25 \\
Coluna errada                 & 19  & 9  & 10 & 0  \\
Outro erro semântico          & 10  & 5  & 5  & 0  \\
Agregação errada              & 5   & 5  & 0  & 0  \\
\bottomrule
\end{tabular}
\end{table}

Como o parser oficial do Spider já se mostrou limitado a uma gramática SQL
restrita neste trabalho (cf. o caso de parênteses agrupando condições no
\texttt{WHERE}, seção anterior), investigou-se se parte dessas 179 previsões
``inválidas'' seriam, na verdade, SQLite sintaticamente válido que o parser
simplesmente não reconhece. Cada uma das 179 previsões foi executada
diretamente contra o banco SQLite real correspondente (\texttt{sqlite3},
sem passar pelo parser AST do Spider). O resultado foi conclusivo:
\textbf{nenhuma das 179 (0\%) executou com sucesso} -- todas produziram um
erro real do SQLite, predominantemente \texttt{no such table} ou
\texttt{no such column} (por exemplo, \texttt{institutions.uncs\_id} em vez
de \texttt{unics\_id}, ou um join com a tabela inexistente \texttt{subjects}
em vez de \texttt{subject\_areas}). Isso descarta a hipótese de artefato de
parser para esta categoria: a maioria dos erros do modelo no
ScienceBenchmark é, de fato, alucinação de nomes de tabela/coluna
inexistentes no esquema ou SQL genuinamente malformado, não uma limitação da
ferramenta de avaliação.

\subsection{Decodificação restrita por gramática (PICARD)}
\label{sec:picard-attempt}

Dado que a maioria dos erros é SQL sintaticamente inválido, testou-se se a
decodificação restrita por gramática do PICARD \cite{picard2021}
-- que impede o modelo de gerar tokens que levariam a uma consulta
inválida -- reduziria essa taxa. O eval foi refeito com
\texttt{use\_picard=true}, \texttt{launch\_picard=true}, mantendo os demais
parâmetros idênticos.

\textbf{Resultado}: as 293 previsões geradas com PICARD ativo são
\emph{idênticas, token a token}, às geradas sem PICARD -- exact match e exec
saíram exatamente iguais aos valores da Tabela~\ref{tab:sciencebenchmark-zeroshot}.
Investigando a causa, constatou-se que o PICARD nunca chega a interceptar a
geração nesta configuração: em \texttt{seq2seq/models/graphix/rgat\_picard.py},
um ajuste anterior (necessário para corrigir a resolução da classe do modelo,
que sem ele carregava um T5 genérico em vez da variante com RGAT) passou a
importar a classe do modelo diretamente, contornando o mecanismo
(\texttt{model\_cls\_wrapper}) pelo qual o \texttt{generate()} restrito do
PICARD normalmente seria aplicado. O servidor PICARD chega a subir e ficar
ativo durante o eval, mas nunca é de fato consultado pela geração --
confirmado tanto pela ausência total de qualquer menção a
``picard''/``thrift'' nos logs do eval quanto pela identidade exata das
previsões. Trata-se de uma lacuna de integração pré-existente, não avaliada
anteriormente porque o eval com PICARD nunca havia sido exercitado de ponta
a ponta neste projeto; o conserto exigiria combinar os dois mecanismos de
substituição de classe (RGAT-aware e PICARD-aware) em vez de tratá-los como
alternativas mutuamente exclusivas, o que fica registrado aqui como limitação
conhecida e trabalho futuro.

\subsection{Busca em feixe (\emph{beam search}, num\_beams=4)}
\label{sec:beam4}

Como variação de decodificação adicional (mantendo o checkpoint zero-shot
inalterado), o eval foi refeito com \texttt{num\_beams=4} no lugar da
decodificação gulosa (\texttt{num\_beams=1}) usada nos resultados acima.

\textbf{Resultado}: 9,34\% exact match / 9,56\% exec (289 exemplos
pontuados para exact match, 287 para exec), \texttt{eval\_loss=0,6885} --
uma melhora pequena mas real sobre a decodificação gulosa (9,0\%/8,87\%),
consistente com o ganho usual, porém modesto, de busca em feixe sobre
geração gulosa em tarefas de geração estruturada.

\section{Fine-tuning no ScienceBenchmark}
\label{sec:finetuning}

A partir do mesmo checkpoint zero-shot, foram conduzidos dois experimentos de
ajuste fino sobre os dados de treino do ScienceBenchmark (4732 exemplos,
seed+synth combinados): um fine-tuning completo (todas as épocas sobre o
maior subconjunto de dados que a GPU disponível comporta) e um fine-tuning
\emph{few-shot} (poucas centenas de exemplos, mais épocas). Ambos foram
avaliados no mesmo split de dev usado no zero-shot, seguindo a mesma
metodologia de scoring (exact match + execução, mesmos filtros de
segurança contra travamentos de execução no \texttt{sdss}).

\textbf{Restrição de hardware determinante}: os experimentos rodaram numa
única GPU de consumidor (NVIDIA GeForce GTX 1070, 8\,GB de VRAM, sem tensor
cores), o mesmo hardware usado para os evals acima. Diferente do Spider
(esquemas pequenos, $\sim$5,3 tabelas/banco em média), os grafos RGAT do
ScienceBenchmark chegam a mais de 1300 nós por exemplo, o que por si só é
suficiente para saturar 8\,GB de VRAM durante o treino (fase que aloca
ativações para retropropagação, ao contrário do eval, que só faz
inferência).

\subsection{Fine-tuning completo}
\label{sec:finetuning-full}

\textbf{Configuração}: 3 épocas, \texttt{batch\_size=1} com
\texttt{gradient\_accumulation\_steps=32} (lote efetivo 32), taxa de
aprendizado $5\times10^{-5}$ com \texttt{adafactor}, 330 passos de
otimização totais. Por restrição de VRAM, exemplos cujo grafo RGAT excede
1000 nós foram descartados do treino -- isso excluiu \textbf{1165 dos 1165
exemplos de \texttt{oncomx}} presentes no split de treino após o primeiro
filtro (praticamente o banco inteiro; ver limitação (iii) abaixo), mantendo
\texttt{cordis} e \texttt{sdss} quase integralmente (3540 dos 4732 exemplos
totais permaneceram).

\begin{table}[h]
\centering
\caption{Fine-tuning completo: resultado final no dev set completo (289--293 exemplos, sem o filtro de tamanho de grafo aplicado ao treino).}
\label{tab:finetuning-full}
\begin{tabular}{lrrr}
\toprule
\textbf{Configuração} & \textbf{exact match} & \textbf{exec} & \textbf{eval\_loss} \\
\midrule
Zero-shot (referência) & 9,0\%  & 8,87\% & -- \\
Fine-tuning completo    & \textbf{0,0\%} & \textbf{0,34\%} & 4,2321 \\
\bottomrule
\end{tabular}
\end{table}

O resultado é \textbf{pior que o zero-shot}, apesar do \emph{loss} de
treino ter caído de forma saudável e monotônica até 0,469 ao final das 3
épocas. A causa raiz identificada: o modelo nunca viu nenhum exemplo de
\texttt{oncomx} durante o treino (excluído por VRAM), mas o dev set de
avaliação final inclui \texttt{oncomx} normalmente -- o desempenho ruim
nesse domínio nunca visto arrasta a média agregada para baixo, revertendo
o ganho que o fine-tuning trouxe em \texttt{cordis}/\texttt{sdss}. Esse é
um achado metodológico central deste trabalho: \textbf{fine-tuning
completo sob restrição severa de VRAM que exclui um dos domínios-alvo pode
piorar a generalização agregada}, mesmo quando todos os indicadores de
treino (perda) sugerem que o ajuste está funcionando.

\subsection{Fine-tuning \emph{few-shot}}
\label{sec:finetuning-fewshot}

Para tentar obter alguma representação de \texttt{oncomx} sem incorrer no
mesmo custo proibitivo de VRAM ao longo de um treino completo, foi montado
um segundo experimento com uma amostra pequena e deliberadamente balanceada
entre os três bancos: 140 exemplos de \texttt{cordis} + 140 de
\texttt{sdss} (amostragem aleatória) + 20 de \texttt{oncomx} (os 20 exemplos
de menor grafo disponíveis, 1022--1025 nós -- ver limitação (iii)), 300
exemplos no total, por 10 épocas (90 passos de otimização).

\begin{table}[h]
\centering
\caption{Fine-tuning few-shot: resultado final no dev set completo.}
\label{tab:finetuning-fewshot}
\begin{tabular}{lrrr}
\toprule
\textbf{Configuração} & \textbf{exact match} & \textbf{exec} & \textbf{eval\_loss} \\
\midrule
Zero-shot (referência)  & 9,0\%  & 8,87\% & -- \\
Fine-tuning completo     & 0,0\%  & 0,34\% & 4,2321 \\
Fine-tuning few-shot     & \textbf{0,0\%} & \textbf{0,0\%} & 4,8754 \\
\bottomrule
\end{tabular}
\end{table}

O resultado também ficou abaixo do zero-shot, e abaixo até do fine-tuning
completo no exec. Diferente do fine-tuning completo, aqui não há exclusão
sistemática de um banco (os três estão representados, ainda que
\texttt{oncomx} de forma mínima) -- a causa mais provável é volume de
dados insuficiente (300 exemplos) combinado com uma limitação metodológica
própria deste experimento, descrita em (iv) abaixo: o checkpoint final
salvo corresponde a apenas 52\% do treino planejado, não ao ponto de
convergência real.

\subsection{Correção do \emph{gradient checkpointing} e fine-tuning completo revisado}
\label{sec:gc-fix}

A limitação (iii) acima -- exclusão total de \texttt{oncomx} do treino por
falta de VRAM -- motivou uma segunda iteração deste trabalho, buscando
incluir os três domínios sem reduzir \texttt{max\_source\_length}. O
mecanismo natural para isso é o \emph{gradient checkpointing} do PyTorch
(\texttt{torch.utils.checkpoint}): em vez de manter em memória todas as
ativações intermediárias da rede para a retropropagação, ele as recalcula
sob demanda durante o \emph{backward pass}, trocando tempo de computação por
memória. A flag já existia em \texttt{configs/train.json}, mas estava sem
efeito para o modelo com RGAT -- dois problemas distintos, ambos
corrigidos:

\begin{itemize}
\item A classe \texttt{Model} em \texttt{seq2seq/models/graphix/rgat.py} é
um invólucro (\emph{wrapper}) que guarda o T5 com RGAT de fato como um
atributo interno (\texttt{self.pretrain\_model}), em vez de herdar sua
hierarquia de camadas. O método padrão
\texttt{gradient\_checkpointing\_enable()} do HuggingFace percorre
recursivamente as submódulos de uma classe para ativar o checkpointing em
cada um -- como o RGAT vive dentro de um atributo e não na hierarquia
percorrida, a ativação nunca chegava lá. Corrigido delegando
explicitamente a chamada ao \texttt{pretrain\_model} interno.
\item Mesmo com a delegação corrigida, o \texttt{T5Stack.forward()} (em
\texttt{seq2seq/models/modeling\_t5.py}) constrói, quando o checkpointing
está ativo, uma função de encapsulamento (\emph{closure}) para cada bloco
que é passada a \texttt{torch.utils.checkpoint.checkpoint()} -- e essa
closure não repassava \texttt{graph\_batch}/\texttt{relation\_emb} ao
bloco, fazendo o RGAT ser silenciosamente ignorado (ou falhar) sempre que
o checkpointing estava ativo. Corrigido capturando esses dois objetos por
\emph{closure} do Python (não são tensores diferenciáveis que o mecanismo
de checkpoint precise gerenciar, então isso é seguro tanto na chamada
inicial quanto no recálculo do \emph{backward}).
\end{itemize}

A correção foi validada com um teste de fumaça isolado sobre os seis
maiores grafos de \texttt{oncomx} (1217--1308 nós, os que mais
frequentemente causavam OOM antes), confirmando treino sem erros e perda
decrescente normalmente.

Com o \emph{gradient checkpointing} funcionando, o fine-tuning completo foi
refeito com os mesmos hiperparâmetros da Seção~\ref{sec:finetuning-full}
(3 épocas, \texttt{batch\_size=1}, \texttt{gradient\_accumulation\_steps=32},
$5\times10^{-5}$ com \texttt{adafactor}), desta vez sem nenhuma exclusão de
exemplo por tamanho de grafo -- os 4732 exemplos de treino (incluindo
\texttt{oncomx} por completo) participaram do treino.

\begin{table}[h]
\centering
\caption{Fine-tuning completo, antes e depois da correção do gradient checkpointing.}
\label{tab:finetuning-full-gc}
\begin{tabular}{lrrrl}
\toprule
\textbf{Configuração} & \textbf{exact match} & \textbf{exec} & \textbf{eval\_loss} & \textbf{train\_loss / tempo} \\
\midrule
Fine-tuning completo (sem \texttt{oncomx}) & 0,0\% & 0,34\% & 4,2321 & 0,469 / -- \\
Fine-tuning completo (com \texttt{oncomx}, gc) & 0,0\% & \textbf{0,34\%} & \textbf{1,6765} & 2,215 / 9h42min \\
\bottomrule
\end{tabular}
\end{table}

O resultado é um achado genuinamente nuançado: incluir \texttt{oncomx}
reduziu o \texttt{eval\_loss} em quase $2{,}5\times$ (4,2321 $\to$ 1,6765) --
uma melhora clara e esperada, já que o modelo agora vê exemplos do domínio
que antes lhe era totalmente estranho. Mas essa melhora \textbf{não se
traduziu em nenhum ganho nas métricas discretas}: \texttt{exact\_match}
permanece em 0\% e \texttt{exec} praticamente inalterado (0,34\%). Isso
sugere que a hipótese da Seção~\ref{sec:finetuning-full} (exclusão de
\texttt{oncomx} como causa única do resultado ruim) era \emph{necessária
mas não suficiente}: mesmo com todos os três domínios representados no
treino, os erros do modelo continuam sendo do tipo que quebra
\texttt{exact\_match}/\texttt{exec} de forma binária -- nome de
tabela/coluna incorreto, estrutura de \texttt{JOIN} equivocada -- e uma
perda agregada mais baixa (que mede o quão "próxima" a distribuição de
tokens gerada está do alvo, token a token) não impede esse tipo de erro
categórico. As causas mais prováveis são um orçamento de treino ainda
pequeno (3 épocas sobre 4732 exemplos heterogêneos, contra 5 épocas sobre
8577 exemplos homogêneos no Spider original) e/ou uma limitação de
capacidade do próprio t5-base diante da diversidade de esquemas do
ScienceBenchmark, discutida em conjunto com o achado da próxima seção.

\subsection{Baseline t5-small sem RGAT}
\label{sec:t5small-baseline}

Para isolar o efeito específico do RGAT dos demais fatores (tamanho do
modelo, orçamento de treino, dados), foi conduzido um experimento de
ablação: um T5 padrão, \textbf{sem nenhuma camada RGAT ou estrutura de
grafo}, na variante \texttt{t5-small} (60M parâmetros, contra os $\sim$220M
do \texttt{t5-base} usado em todos os experimentos anteriores). O restante
do pipeline foi mantido idêntico -- mesmo pré-processamento (a mesma
serialização de esquema textual produzida pelas etapas de grafo, ainda que
o grafo em si não seja consumido), mesmo split de treino/dev (4732/299
exemplos), mesmos hiperparâmetros (3 épocas, \texttt{batch\_size=1},
\texttt{gradient\_accumulation\_steps=32}, $5\times10^{-5}$ com
\texttt{adafactor}). Uma nova classe \texttt{Model}
(\texttt{seq2seq/models/graphix/plain.py}) carrega um
\texttt{T5ForConditionalGeneration} padrão do HuggingFace e ignora
completamente o \texttt{graph\_batch}, selecionável por uma variável de
ambiente (\texttt{GRAPHIX\_MODEL\_VARIANT=plain}) sem duplicar o restante
do script de treino.

\begin{table}[h]
\centering
\caption{t5-small sem RGAT vs.\ t5-base com RGAT (fine-tuning completo).}
\label{tab:t5small-vs-rgat}
\begin{tabular}{lrrrl}
\toprule
\textbf{Configuração} & \textbf{exact match} & \textbf{exec} & \textbf{eval\_loss} & \textbf{tempo de treino} \\
\midrule
t5-base + RGAT (com \texttt{oncomx}, gc) & 0,0\% & 0,34\% & 1,6765 & 9h42min \\
t5-small \textbf{sem} RGAT (3 épocas)     & \textbf{1,38\%} & \textbf{1,37\%} & \textbf{1,1825} & \textbf{2h53min} \\
\bottomrule
\end{tabular}
\end{table}

O resultado é surpreendente: o modelo \emph{menor} e \emph{sem} qualquer
estrutura de grafo superou o modelo com RGAT em todas as métricas -- exact
match 4$\times$ maior (1,38\% vs.\ 0,34\% de exec no RGAT, comparação
direta de exec: 1,37\% vs.\ 0,34\%, também $\sim$4$\times$), \texttt{eval\_loss}
menor, e treino $3{,}4\times$ mais rápido (sem o custo de recomputação do
\emph{gradient checkpointing} nem a sobrecarga de propagar mensagens em
grafos de até 1300 nós por exemplo). É importante notar que essa
comparação \textbf{não é perfeitamente controlada}: t5-small e t5-base+RGAT
diferem tanto na presença do RGAT quanto no tamanho do modelo base, então
o experimento isola "RGAT vs.\ sem RGAT" apenas de forma aproximada, não
com uma variável por vez. Ainda assim, o resultado é informativo na direção
oposta à esperada -- se o RGAT estivesse agregando valor líquido, seria
razoável esperar que o modelo \emph{maior e com estrutura adicional}
superasse o menor e mais simples, não o contrário. Hipóteses para essa
inversão:

\begin{itemize}
\item O RGAT deste projeto foi desenhado e ajustado para os esquemas
pequenos e homogêneos do Spider ($\sim$5,3 tabelas/banco); grafos de
1000+ nós do ScienceBenchmark estão 2--3 ordens de grandeza acima da escala
para a qual a arquitetura foi validada, e a propagação de mensagens em
grafos tão densos pode diluir sinal útil em vez de reforçá-lo.
\item Um modelo maior e mais complexo (t5-base+RGAT) tipicamente precisa de
mais dados e/ou mais épocas para superar um modelo menor em fine-tuning --
com apenas 4732 exemplos e 3 épocas, o orçamento de treino pode ser
insuficiente para que a capacidade extra do RGAT se traduza em vantagem
prática, enquanto o t5-small, por ser mais simples, converge mais
rapidamente para um padrão razoável com o mesmo orçamento.
\item A complexidade adicional (mais parâmetros, mais uma fonte de
gradiente via o \emph{encoder} do grafo) pode estar introduzindo ruído de
otimização que atrapalha a convergência dentro do orçamento de treino
disponível, sem chegar a comprometer a estabilidade geral (o \emph{loss}
de treino do RGAT+t5-base ainda caiu de forma saudável, só não tão baixo
quanto seria necessário para competir com o t5-small nas métricas
discretas).
\end{itemize}

Este resultado não permite concluir que RGAT seja inerentemente inadequado
para text-to-SQL sobre esquemas grandes -- apenas que, \textbf{nas
condições específicas deste trabalho} (hardware de 8\,GB de VRAM, poucas
horas de treino disponíveis, 4732 exemplos), a complexidade adicional do
RGAT não se pagou em métricas discretas frente a um baseline T5 puro e
menor, treinado com o mesmo orçamento de tempo.

\subsection{Limitações metodológicas}
\label{sec:limitations}

Os dois experimentos de fine-tuning expuseram uma série de limitações,
tanto de hardware quanto de configuração, documentadas aqui por
transparência e como recomendação para trabalho futuro:

\begin{enumerate}
\item[(i)] \textbf{Ausência de split de teste genuíno.} O ScienceBenchmark,
como distribuído, só define \texttt{train} (seed+synth) e \texttt{dev}. O
mesmo split de \texttt{dev} usado para validação durante o treino
(seleção do "melhor" checkpoint) foi reutilizado como avaliação final --
não há um terceiro split, cego durante todo o processo, para reportar como
número de teste independente. O Spider, usado no treino-base, tem a mesma
limitação estrutural (sem test set público), então isso é consistente com
a metodologia herdada, mas permanece uma ressalva válida: o número final
reportado pode estar levemente inflado em relação a um teste verdadeiramente
cego.

\item[(ii)] \textbf{\texttt{torch.cuda.empty\_cache()} incondicional.} Uma
chamada incondicional a \texttt{torch.cuda.empty\_cache()} em
\texttt{seq2seq/models/graphix/rgat\_tuning.py::RGAT\_Layer.forward}, uma
vez por camada RGAT (12$\times$ por exemplo no t5-base), forçava uma
sincronização CUDA completa a cada chamada. Esse código é compartilhado
entre treino e avaliação (\texttt{modeling\_t5.py} importa a mesma classe
para ambos), então os resultados de zero-shot/PICARD/beam4 relatados acima
provavelmente rodaram mais devagar do que precisavam -- os números
continuam válidos, só o tempo de execução foi maior que o necessário. O
fix (só chamar \texttt{empty\_cache()} sob OOM real, não preventivamente)
foi validado por comparação direta: antes, nenhum passo de otimização
completava em mais de uma hora; depois, o mesmo passo completava em
minutos.

\item[(iii)] \textbf{\texttt{oncomx} tem esquema uniformemente grande
demais para 8\,GB de VRAM.} Todos os 1165 exemplos de treino de
\texttt{oncomx} produzem grafos RGAT entre 1022 e 1308 nós -- não existe
um subconjunto pequeno e seguro para amostrar. Mesmo os 20 exemplos
"menores" usados no few-shot (1022--1025 nós) ainda estão acima do teto de
900--1000 nós que se mostrou confiável durante o treino completo, e
provocaram OOMs individuais recorrentes durante o few-shot (capturados e
pulados por um mecanismo de segurança, não por acaso -- ver item v). Esta é
uma restrição de hardware, não uma limitação de amostragem: nenhuma
estratégia de seleção de exemplos resolve o problema sem reduzir
\texttt{max\_source\_length} (o que exigiria reprocessar toda a base de
grafos RGAT, fora do escopo deste trabalho). \textbf{Atualização}: esta
restrição foi contornada numa iteração posterior deste trabalho habilitando
corretamente o \emph{gradient checkpointing} do PyTorch, o que permitiu
incluir \texttt{oncomx} por completo no treino sem reduzir
\texttt{max\_source\_length} -- ver Seção~\ref{sec:gc-fix}.

\item[(iv)] \textbf{Ausência de \emph{warmup} de taxa de aprendizado no
fine-tuning completo, e espaçamento de checkpoint excessivo no few-shot.}
O fine-tuning completo usou \texttt{warmup\_ratio=0}, fazendo a taxa de
aprendizado saltar diretamente para $5\times10^{-5}$ no primeiro passo --
isso corresponde a um colapso temporário de geração (repetição de
n-gramas, ex.\ \texttt{"...loadversion = 301 and loadversion = 301 and..."}
até o limite de tokens), causando \texttt{eval\_exact\_match=0} em todos os
checkpoints intermediários mesmo com o \emph{loss} caindo normalmente. Um
\texttt{no\_repeat\_ngram\_size=3} foi adicionado durante o treino
(reduzindo o tempo de geração em $5\times$ e eliminando os loops de
n-grama), e o few-shot já nasceu com \texttt{warmup\_ratio=0,1} e
\texttt{no\_repeat\_ngram\_size=3} desde o início -- mas um padrão de
repetição mais sutil (alternância entre tokens SentencePiece de tamanhos
diferentes que decodificam para o mesmo caractere repetido, ex.\
\texttt{"6666666..."}) persistiu, sugerindo que o modelo continuou
sub-treinado para geração de literais numéricos mesmo ao final do processo.
Separadamente, no few-shot, o intervalo entre checkpoints
(\texttt{eval\_steps=47}) foi ajustado no meio do treino para reduzir o
tempo gasto em avaliação (cada rodada de eval levava até 111 minutos), mas
isso fez com que só um checkpoint intermediário (passo 47 de 90, 52\% do
treino) fosse salvo antes do fim -- o \texttt{load\_best\_model\_at\_end}
carregou esse checkpoint como "melhor" por ser o único candidato com
métrica registrada, e os pesos do passo 90 (o treino de fato mais avançado)
nunca chegaram a ser persistidos em disco.

\item[(v)] \textbf{Seleção de "melhor checkpoint" quebrada quando a métrica
empata.} O fine-tuning completo usou
\texttt{metric\_for\_best\_model="exact\_match"}. Como o colapso de geração
descrito em (iv) manteve \texttt{exact\_match=0,0} em \emph{todos} os
checkpoints intermediários, e a lógica de comparação do HuggingFace
Trainer só substitui o "melhor" registrado quando a nova métrica é
\emph{estritamente} maior, o primeiro checkpoint avaliado (passo 50 de
330, a versão \emph{menos} treinada de todas) nunca foi superado e acabou
sendo o modelo efetivamente salvo ao final do treino -- os pesos do passo
300 (97\% do treino) precisaram ser recuperados manualmente do diretório de
checkpoint correspondente para a avaliação real reportada na
Tabela~\ref{tab:finetuning-full}. O few-shot corrigiu isso usando
\texttt{metric\_for\_best\_model="eval\_loss"} com
\texttt{greater\_is\_better=false} (uma métrica contínua, que decresceu de
forma monotônica ao longo do treino), mas essa correção por si só não
evitou a limitação (iv) relacionada ao espaçamento de checkpoints.

\item[(vi)] \textbf{OOMs genuínos e travamento de driver de GPU.} Dois
crashes de OOM (\emph{out of memory}) genuínos interromperam o fine-tuning
completo no meio do treino, e um deles coincidiu com um travamento
completo do driver NVIDIA (\texttt{nvidia-smi} parou de reconhecer a
placa), exigindo reinicialização completa da máquina -- um reset de driver
isolado não foi suficiente para restaurar o desempenho normal da GPU.
Mitigações aplicadas: \emph{checkpointing} a cada 50 passos (em vez de 500,
o intervalo original herdado do treino no Spider, que nunca dispararia
dentro de um treino de 330 passos) e uma captura de OOM por micro-batch em
\texttt{seq2seq/utils/trainer.py::Seq2SeqTrainer.training\_step} que pula
apenas o exemplo problemático (zerando o gradiente parcialmente acumulado
e liberando cache) em vez de derrubar o processo inteiro -- validado
diretamente durante o few-shot, que sobreviveu a dezenas de OOMs
individuais de exemplos grandes de \texttt{oncomx} sem interrupção.
\end{enumerate}

\section{Comparação final}
\label{sec:final-comparison}

\begin{table}[h]
\centering
\caption{Comparação de todas as configurações avaliadas no dev set do ScienceBenchmark.}
\label{tab:final-comparison}
\begin{tabular}{lrrl}
\toprule
\textbf{Configuração} & \textbf{exact match} & \textbf{exec} & \textbf{Observação} \\
\midrule
Zero-shot (greedy, num\_beams=1)  & 9,0\%  & 8,87\% & referência \\
Zero-shot + beam search (n=4)      & \textbf{9,34\%} & \textbf{9,56\%} & melhor resultado geral \\
Zero-shot + PICARD                 & 9,0\%  & 8,87\% & PICARD inerte (falha de integração, \S\ref{sec:picard-attempt}) \\
Fine-tuning completo, sem \texttt{oncomx} (3 épocas)    & 0,0\%  & 0,34\% & pior que zero-shot (\S\ref{sec:finetuning-full}) \\
Fine-tuning few-shot (300 ex., 10 épocas) & 0,0\% & 0,0\% & pior que zero-shot (\S\ref{sec:finetuning-fewshot}) \\
Fine-tuning completo, com \texttt{oncomx} + gc (3 épocas) & 0,0\% & 0,34\% & eval\_loss $2,5\times$ menor, métricas discretas inalteradas (\S\ref{sec:gc-fix}) \\
Baseline t5-small sem RGAT (3 épocas)     & 1,38\% & 1,37\% & supera todos os fine-tunings com RGAT (\S\ref{sec:t5small-baseline}) \\
\bottomrule
\end{tabular}
\end{table}

O achado central deste conjunto de experimentos permanece, em essência,
\textbf{negativo em relação ao fine-tuning com RGAT}: em nenhuma das quatro
tentativas de ajuste fino (completo sem \texttt{oncomx}, \emph{few-shot},
completo com \texttt{oncomx} via \emph{gradient checkpointing}) o modelo
t5-base+RGAT superou o desempenho do checkpoint zero-shot com busca em
feixe nas métricas discretas -- mesmo depois de resolver a limitação de
VRAM que impedia incluir todos os domínios de treino. A correção do
\emph{gradient checkpointing} (\S\ref{sec:gc-fix}) trouxe uma melhora real
e mensurável na perda de avaliação (2,5$\times$ menor), mas essa melhora
não se propagou às métricas discretas -- um resultado que, por si só, já é
instrutivo: perda agregada e acerto exato/execução podem divergir
significativamente quando os erros do modelo são predominantemente
categóricos (nome de tabela/coluna errado) em vez de graduais.

O achado mais forte deste trabalho, no entanto, veio do baseline de
ablação (\S\ref{sec:t5small-baseline}): um \texttt{t5-small} \textbf{sem
nenhuma estrutura de grafo}, treinado com o mesmo orçamento de dados e
épocas, superou \emph{todas} as configurações com RGAT fine-tunadas neste
trabalho em exact match e exec, e o fez em menos de um terço do tempo de
treino. Combinado com o resultado do zero-shot (onde o RGAT+t5-base é a
única opção avaliada, pois não há um t5-small equivalente pré-treinado no
Spider para comparação direta), a evidência deste trabalho aponta para uma
conclusão qualificada: \textbf{nas condições de hardware e dados
disponíveis, a complexidade adicional do RGAT não demonstrou vantagem
sobre um T5 padrão no fine-tuning para o ScienceBenchmark} -- um resultado
negativo genuíno para a hipótese de que a estrutura de grafo ajudaria a
generalizar para esquemas maiores e de domínio distinto do treino
original. As limitações documentadas na Seção~\ref{sec:limitations}
(volume de dados, orçamento de épocas, ausência de um terceiro split de
teste cego) continuam válidas e sugerem que essa conclusão é específica
deste orçamento experimental, não necessariamente uma propriedade
fundamental da arquitetura RGAT -- mas dentro do escopo efetivamente
testado, o T5 puro foi tanto mais eficaz quanto mais barato de treinar.

% ou compilado standalone.

\section{Avaliação no ScienceBenchmark (zero-shot)}
\label{sec:sciencebenchmark-eval}

O checkpoint treinado no Spider (t5-base + RGAT, 5 épocas, 8577 exemplos de
treino, resultado no dev set do Spider: 41,9\% \emph{exact match} / 44,8\%
\emph{execution accuracy}) foi avaliado, sem nenhum re-treinamento, sobre os
299 exemplos do split de desenvolvimento do ScienceBenchmark
\cite{sciencebenchmark2023}, distribuídos em três bancos de domínios distintos:
\texttt{cordis} (projetos de pesquisa europeus), \texttt{oncomx} (biomedicina)
e \texttt{sdss} (astronomia).

\subsection{Resultado agregado}

\begin{table}[h]
\centering
\caption{Resultado do eval zero-shot do checkpoint Spider-treinado no ScienceBenchmark, por banco.}
\label{tab:sciencebenchmark-zeroshot}
\begin{tabular}{lrrrl}
\toprule
\textbf{Banco} & \textbf{n} & \textbf{exact match} & \textbf{exec} & \textbf{Exemplos avaliados} \\
\midrule
cordis            & 100 & 12,0\% & 15,0\% & 100/100 \\
oncomx            & 99  & 14,1\% & 11,1\% & 99/99 (exec: 97/99) \\
sdss              & 94  & 0,0\%  & 0,0\%  & 90/94\textsuperscript{a} \\
\midrule
\textbf{Total (agregado)} & \textbf{293} & \textbf{9,0\%} & \textbf{8,9\%} & exact: 289/293, exec: 287/293 \\
\bottomrule
\end{tabular}

\vspace{0.3em}
{\footnotesize\textsuperscript{a}4 exemplos do sdss descartados da métrica: parser
SQL do Spider (\texttt{third\_party/spider/process\_sql.py}) não suporta
parênteses agrupando condições no \texttt{WHERE}, um padrão ausente no SQL
gold do Spider mas presente no do ScienceBenchmark.}
\end{table}

Para efeito de comparação, o mesmo checkpoint avaliado no dev set do Spider
(1034 exemplos, mesmo domínio de treino) obteve 41,9\% exact match e 44,8\%
exec -- uma queda de $\sim$33 pontos percentuais ao mudar de domínio sem
re-treinamento, consistente com a expectativa de que arquiteturas
graph-aware treinadas em um conjunto de bancos pequenos e homogêneos (Spider,
$\sim$5,3 tabelas/banco em média) generalizam de forma limitada para bancos
maiores e de domínio mais específico (ScienceBenchmark, 19--25 tabelas no
\texttt{cordis}, colunas de medição numérica pouco usuais no \texttt{sdss}).

\subsection{Quebra por dificuldade}

Usando a classificação de dificuldade do avaliador oficial do Spider
(\emph{easy}/\emph{medium}/\emph{hard}/\emph{extra}, baseada em número de
componentes SQL, condições e agregações), o exact match cai a zero em
\emph{hard} e \emph{extra} nos três bancos, sem exceção:

\begin{table}[h]
\centering
\caption{Exact match por nível de dificuldade e banco.}
\label{tab:sciencebenchmark-hardness}
\begin{tabular}{lrrrr}
\toprule
\textbf{Banco} & \textbf{easy} & \textbf{medium} & \textbf{hard} & \textbf{extra} \\
\midrule
cordis (n=100) & 32,0\% (n=25) & 11,4\% (n=35) & 0\% (n=19) & 0\% (n=21) \\
oncomx (n=99)  & 42,9\% (n=21) & 13,9\% (n=36) & 0\% (n=22) & 0\% (n=20) \\
sdss (n=90)    & 0\% (n=6)     & 0\% (n=28)    & 0\% (n=20) & 0\% (n=36) \\
\bottomrule
\end{tabular}
\end{table}

O padrão observado sugere duas causas de queda distintas e sobrepostas: (i)
uma limitação de \emph{capacidade}, comum aos três bancos -- o modelo t5-base
resolve consultas simples mas não lida com consultas complexas (múltiplos
JOINs, subconsultas, agregações compostas) fora do domínio de treino; e (ii)
uma limitação de \emph{transferência de domínio}, específica do
\texttt{sdss} -- o resultado é zero mesmo nas consultas \emph{easy}, o que não
ocorre em \texttt{cordis}/\texttt{oncomx}, indicando que o vocabulário e a
estrutura de esquema do \texttt{sdss} (medições fotométricas, nomenclatura
astronômica) são suficientemente distantes do domínio de treino a ponto de
comprometer até os casos mais simples.

Uma observação metodológica reforça essa segunda leitura: o tamanho do
esquema (contagem de nós do grafo RGAT) e o banco de dados são quase
colineares nesta amostra -- \texttt{cordis}/\texttt{oncomx} sempre excedem
512 nós, \texttt{sdss} está sempre abaixo. Como \texttt{sdss} é o banco com o
\emph{menor} esquema em número de nós e ainda assim o único com resultado
nulo em todos os níveis de dificuldade, o fator determinante da queda não
parece ser o tamanho do esquema, e sim a distância de domínio/vocabulário em
relação ao Spider.

\subsection{Análise qualitativa dos erros}

Das 263 previsões incorretas (excluindo os 4 exemplos de \texttt{sdss} sem
gold parseável), 179 (68\%) foram inicialmente classificadas como ``SQL
sintaticamente inválido'' pelo parser oficial do Spider
(\texttt{third\_party/spider/process\_sql.py}) -- ou seja, o modelo nem
produziu uma consulta que o parser conseguisse interpretar como AST SQL.

\begin{table}[h]
\centering
\caption{Distribuição das categorias de erro (previsões incorretas, exact match = 0).}
\label{tab:error-categories}
\begin{tabular}{lrrrr}
\toprule
\textbf{Categoria} & \textbf{Total} & \textbf{cordis} & \textbf{oncomx} & \textbf{sdss} \\
\midrule
SQL sintaticamente inválido & 179 & 56 & 58 & 65 \\
Tabela errada                & 50  & 13 & 12 & 25 \\
Coluna errada                 & 19  & 9  & 10 & 0  \\
Outro erro semântico          & 10  & 5  & 5  & 0  \\
Agregação errada              & 5   & 5  & 0  & 0  \\
\bottomrule
\end{tabular}
\end{table}

Como o parser oficial do Spider já se mostrou limitado a uma gramática SQL
restrita neste trabalho (cf. o caso de parênteses agrupando condições no
\texttt{WHERE}, seção anterior), investigou-se se parte dessas 179 previsões
``inválidas'' seriam, na verdade, SQLite sintaticamente válido que o parser
simplesmente não reconhece. Cada uma das 179 previsões foi executada
diretamente contra o banco SQLite real correspondente (\texttt{sqlite3},
sem passar pelo parser AST do Spider). O resultado foi conclusivo:
\textbf{nenhuma das 179 (0\%) executou com sucesso} -- todas produziram um
erro real do SQLite, predominantemente \texttt{no such table} ou
\texttt{no such column} (por exemplo, \texttt{institutions.uncs\_id} em vez
de \texttt{unics\_id}, ou um join com a tabela inexistente \texttt{subjects}
em vez de \texttt{subject\_areas}). Isso descarta a hipótese de artefato de
parser para esta categoria: a maioria dos erros do modelo no
ScienceBenchmark é, de fato, alucinação de nomes de tabela/coluna
inexistentes no esquema ou SQL genuinamente malformado, não uma limitação da
ferramenta de avaliação.

\subsection{Decodificação restrita por gramática (PICARD)}
\label{sec:picard-attempt}

Dado que a maioria dos erros é SQL sintaticamente inválido, testou-se se a
decodificação restrita por gramática do PICARD \cite{picard2021}
-- que impede o modelo de gerar tokens que levariam a uma consulta
inválida -- reduziria essa taxa. O eval foi refeito com
\texttt{use\_picard=true}, \texttt{launch\_picard=true}, mantendo os demais
parâmetros idênticos.

\textbf{Resultado}: as 293 previsões geradas com PICARD ativo são
\emph{idênticas, token a token}, às geradas sem PICARD -- exact match e exec
saíram exatamente iguais aos valores da Tabela~\ref{tab:sciencebenchmark-zeroshot}.
Investigando a causa, constatou-se que o PICARD nunca chega a interceptar a
geração nesta configuração: em \texttt{seq2seq/models/graphix/rgat\_picard.py},
um ajuste anterior (necessário para corrigir a resolução da classe do modelo,
que sem ele carregava um T5 genérico em vez da variante com RGAT) passou a
importar a classe do modelo diretamente, contornando o mecanismo
(\texttt{model\_cls\_wrapper}) pelo qual o \texttt{generate()} restrito do
PICARD normalmente seria aplicado. O servidor PICARD chega a subir e ficar
ativo durante o eval, mas nunca é de fato consultado pela geração --
confirmado tanto pela ausência total de qualquer menção a
``picard''/``thrift'' nos logs do eval quanto pela identidade exata das
previsões. Trata-se de uma lacuna de integração pré-existente, não avaliada
anteriormente porque o eval com PICARD nunca havia sido exercitado de ponta
a ponta neste projeto; o conserto exigiria combinar os dois mecanismos de
substituição de classe (RGAT-aware e PICARD-aware) em vez de tratá-los como
alternativas mutuamente exclusivas, o que fica registrado aqui como limitação
conhecida e trabalho futuro.

\subsection{Busca em feixe (\emph{beam search}, num\_beams=4)}
\label{sec:beam4}

Como variação de decodificação adicional (mantendo o checkpoint zero-shot
inalterado), o eval foi refeito com \texttt{num\_beams=4} no lugar da
decodificação gulosa (\texttt{num\_beams=1}) usada nos resultados acima.

\textbf{Resultado}: 9,34\% exact match / 9,56\% exec (289 exemplos
pontuados para exact match, 287 para exec), \texttt{eval\_loss=0,6885} --
uma melhora pequena mas real sobre a decodificação gulosa (9,0\%/8,87\%),
consistente com o ganho usual, porém modesto, de busca em feixe sobre
geração gulosa em tarefas de geração estruturada.

\section{Fine-tuning no ScienceBenchmark}
\label{sec:finetuning}

A partir do mesmo checkpoint zero-shot, foram conduzidos dois experimentos de
ajuste fino sobre os dados de treino do ScienceBenchmark (4732 exemplos,
seed+synth combinados): um fine-tuning completo (todas as épocas sobre o
maior subconjunto de dados que a GPU disponível comporta) e um fine-tuning
\emph{few-shot} (poucas centenas de exemplos, mais épocas). Ambos foram
avaliados no mesmo split de dev usado no zero-shot, seguindo a mesma
metodologia de scoring (exact match + execução, mesmos filtros de
segurança contra travamentos de execução no \texttt{sdss}).

\textbf{Restrição de hardware determinante}: os experimentos rodaram numa
única GPU de consumidor (NVIDIA GeForce GTX 1070, 8\,GB de VRAM, sem tensor
cores), o mesmo hardware usado para os evals acima. Diferente do Spider
(esquemas pequenos, $\sim$5,3 tabelas/banco em média), os grafos RGAT do
ScienceBenchmark chegam a mais de 1300 nós por exemplo, o que por si só é
suficiente para saturar 8\,GB de VRAM durante o treino (fase que aloca
ativações para retropropagação, ao contrário do eval, que só faz
inferência).

\subsection{Fine-tuning completo}
\label{sec:finetuning-full}

\textbf{Configuração}: 3 épocas, \texttt{batch\_size=1} com
\texttt{gradient\_accumulation\_steps=32} (lote efetivo 32), taxa de
aprendizado $5\times10^{-5}$ com \texttt{adafactor}, 330 passos de
otimização totais. Por restrição de VRAM, exemplos cujo grafo RGAT excede
1000 nós foram descartados do treino -- isso excluiu \textbf{1165 dos 1165
exemplos de \texttt{oncomx}} presentes no split de treino após o primeiro
filtro (praticamente o banco inteiro; ver limitação (iii) abaixo), mantendo
\texttt{cordis} e \texttt{sdss} quase integralmente (3540 dos 4732 exemplos
totais permaneceram).

\begin{table}[h]
\centering
\caption{Fine-tuning completo: resultado final no dev set completo (289--293 exemplos, sem o filtro de tamanho de grafo aplicado ao treino).}
\label{tab:finetuning-full}
\begin{tabular}{lrrr}
\toprule
\textbf{Configuração} & \textbf{exact match} & \textbf{exec} & \textbf{eval\_loss} \\
\midrule
Zero-shot (referência) & 9,0\%  & 8,87\% & -- \\
Fine-tuning completo    & \textbf{0,0\%} & \textbf{0,34\%} & 4,2321 \\
\bottomrule
\end{tabular}
\end{table}

O resultado é \textbf{pior que o zero-shot}, apesar do \emph{loss} de
treino ter caído de forma saudável e monotônica até 0,469 ao final das 3
épocas. A causa raiz identificada: o modelo nunca viu nenhum exemplo de
\texttt{oncomx} durante o treino (excluído por VRAM), mas o dev set de
avaliação final inclui \texttt{oncomx} normalmente -- o desempenho ruim
nesse domínio nunca visto arrasta a média agregada para baixo, revertendo
o ganho que o fine-tuning trouxe em \texttt{cordis}/\texttt{sdss}. Esse é
um achado metodológico central deste trabalho: \textbf{fine-tuning
completo sob restrição severa de VRAM que exclui um dos domínios-alvo pode
piorar a generalização agregada}, mesmo quando todos os indicadores de
treino (perda) sugerem que o ajuste está funcionando.

\subsection{Fine-tuning \emph{few-shot}}
\label{sec:finetuning-fewshot}

Para tentar obter alguma representação de \texttt{oncomx} sem incorrer no
mesmo custo proibitivo de VRAM ao longo de um treino completo, foi montado
um segundo experimento com uma amostra pequena e deliberadamente balanceada
entre os três bancos: 140 exemplos de \texttt{cordis} + 140 de
\texttt{sdss} (amostragem aleatória) + 20 de \texttt{oncomx} (os 20 exemplos
de menor grafo disponíveis, 1022--1025 nós -- ver limitação (iii)), 300
exemplos no total, por 10 épocas (90 passos de otimização).

\begin{table}[h]
\centering
\caption{Fine-tuning few-shot: resultado final no dev set completo.}
\label{tab:finetuning-fewshot}
\begin{tabular}{lrrr}
\toprule
\textbf{Configuração} & \textbf{exact match} & \textbf{exec} & \textbf{eval\_loss} \\
\midrule
Zero-shot (referência)  & 9,0\%  & 8,87\% & -- \\
Fine-tuning completo     & 0,0\%  & 0,34\% & 4,2321 \\
Fine-tuning few-shot     & \textbf{0,0\%} & \textbf{0,0\%} & 4,8754 \\
\bottomrule
\end{tabular}
\end{table}

O resultado também ficou abaixo do zero-shot, e abaixo até do fine-tuning
completo no exec. Diferente do fine-tuning completo, aqui não há exclusão
sistemática de um banco (os três estão representados, ainda que
\texttt{oncomx} de forma mínima) -- a causa mais provável é volume de
dados insuficiente (300 exemplos) combinado com uma limitação metodológica
própria deste experimento, descrita em (iv) abaixo: o checkpoint final
salvo corresponde a apenas 52\% do treino planejado, não ao ponto de
convergência real.

\subsection{Correção do \emph{gradient checkpointing} e fine-tuning completo revisado}
\label{sec:gc-fix}

A limitação (iii) acima -- exclusão total de \texttt{oncomx} do treino por
falta de VRAM -- motivou uma segunda iteração deste trabalho, buscando
incluir os três domínios sem reduzir \texttt{max\_source\_length}. O
mecanismo natural para isso é o \emph{gradient checkpointing} do PyTorch
(\texttt{torch.utils.checkpoint}): em vez de manter em memória todas as
ativações intermediárias da rede para a retropropagação, ele as recalcula
sob demanda durante o \emph{backward pass}, trocando tempo de computação por
memória. A flag já existia em \texttt{configs/train.json}, mas estava sem
efeito para o modelo com RGAT -- dois problemas distintos, ambos
corrigidos:

\begin{itemize}
\item A classe \texttt{Model} em \texttt{seq2seq/models/graphix/rgat.py} é
um invólucro (\emph{wrapper}) que guarda o T5 com RGAT de fato como um
atributo interno (\texttt{self.pretrain\_model}), em vez de herdar sua
hierarquia de camadas. O método padrão
\texttt{gradient\_checkpointing\_enable()} do HuggingFace percorre
recursivamente as submódulos de uma classe para ativar o checkpointing em
cada um -- como o RGAT vive dentro de um atributo e não na hierarquia
percorrida, a ativação nunca chegava lá. Corrigido delegando
explicitamente a chamada ao \texttt{pretrain\_model} interno.
\item Mesmo com a delegação corrigida, o \texttt{T5Stack.forward()} (em
\texttt{seq2seq/models/modeling\_t5.py}) constrói, quando o checkpointing
está ativo, uma função de encapsulamento (\emph{closure}) para cada bloco
que é passada a \texttt{torch.utils.checkpoint.checkpoint()} -- e essa
closure não repassava \texttt{graph\_batch}/\texttt{relation\_emb} ao
bloco, fazendo o RGAT ser silenciosamente ignorado (ou falhar) sempre que
o checkpointing estava ativo. Corrigido capturando esses dois objetos por
\emph{closure} do Python (não são tensores diferenciáveis que o mecanismo
de checkpoint precise gerenciar, então isso é seguro tanto na chamada
inicial quanto no recálculo do \emph{backward}).
\end{itemize}

A correção foi validada com um teste de fumaça isolado sobre os seis
maiores grafos de \texttt{oncomx} (1217--1308 nós, os que mais
frequentemente causavam OOM antes), confirmando treino sem erros e perda
decrescente normalmente.

Com o \emph{gradient checkpointing} funcionando, o fine-tuning completo foi
refeito com os mesmos hiperparâmetros da Seção~\ref{sec:finetuning-full}
(3 épocas, \texttt{batch\_size=1}, \texttt{gradient\_accumulation\_steps=32},
$5\times10^{-5}$ com \texttt{adafactor}), desta vez sem nenhuma exclusão de
exemplo por tamanho de grafo -- os 4732 exemplos de treino (incluindo
\texttt{oncomx} por completo) participaram do treino.

\begin{table}[h]
\centering
\caption{Fine-tuning completo, antes e depois da correção do gradient checkpointing.}
\label{tab:finetuning-full-gc}
\begin{tabular}{lrrrl}
\toprule
\textbf{Configuração} & \textbf{exact match} & \textbf{exec} & \textbf{eval\_loss} & \textbf{train\_loss / tempo} \\
\midrule
Fine-tuning completo (sem \texttt{oncomx}) & 0,0\% & 0,34\% & 4,2321 & 0,469 / -- \\
Fine-tuning completo (com \texttt{oncomx}, gc) & 0,0\% & \textbf{0,34\%} & \textbf{1,6765} & 2,215 / 9h42min \\
\bottomrule
\end{tabular}
\end{table}

O resultado é um achado genuinamente nuançado: incluir \texttt{oncomx}
reduziu o \texttt{eval\_loss} em quase $2{,}5\times$ (4,2321 $\to$ 1,6765) --
uma melhora clara e esperada, já que o modelo agora vê exemplos do domínio
que antes lhe era totalmente estranho. Mas essa melhora \textbf{não se
traduziu em nenhum ganho nas métricas discretas}: \texttt{exact\_match}
permanece em 0\% e \texttt{exec} praticamente inalterado (0,34\%). Isso
sugere que a hipótese da Seção~\ref{sec:finetuning-full} (exclusão de
\texttt{oncomx} como causa única do resultado ruim) era \emph{necessária
mas não suficiente}: mesmo com todos os três domínios representados no
treino, os erros do modelo continuam sendo do tipo que quebra
\texttt{exact\_match}/\texttt{exec} de forma binária -- nome de
tabela/coluna incorreto, estrutura de \texttt{JOIN} equivocada -- e uma
perda agregada mais baixa (que mede o quão "próxima" a distribuição de
tokens gerada está do alvo, token a token) não impede esse tipo de erro
categórico. As causas mais prováveis são um orçamento de treino ainda
pequeno (3 épocas sobre 4732 exemplos heterogêneos, contra 5 épocas sobre
8577 exemplos homogêneos no Spider original) e/ou uma limitação de
capacidade do próprio t5-base diante da diversidade de esquemas do
ScienceBenchmark, discutida em conjunto com o achado da próxima seção.

\subsection{Baseline t5-small sem RGAT}
\label{sec:t5small-baseline}

Para isolar o efeito específico do RGAT dos demais fatores (tamanho do
modelo, orçamento de treino, dados), foi conduzido um experimento de
ablação: um T5 padrão, \textbf{sem nenhuma camada RGAT ou estrutura de
grafo}, na variante \texttt{t5-small} (60M parâmetros, contra os $\sim$220M
do \texttt{t5-base} usado em todos os experimentos anteriores). O restante
do pipeline foi mantido idêntico -- mesmo pré-processamento (a mesma
serialização de esquema textual produzida pelas etapas de grafo, ainda que
o grafo em si não seja consumido), mesmo split de treino/dev (4732/299
exemplos), mesmos hiperparâmetros (3 épocas, \texttt{batch\_size=1},
\texttt{gradient\_accumulation\_steps=32}, $5\times10^{-5}$ com
\texttt{adafactor}). Uma nova classe \texttt{Model}
(\texttt{seq2seq/models/graphix/plain.py}) carrega um
\texttt{T5ForConditionalGeneration} padrão do HuggingFace e ignora
completamente o \texttt{graph\_batch}, selecionável por uma variável de
ambiente (\texttt{GRAPHIX\_MODEL\_VARIANT=plain}) sem duplicar o restante
do script de treino.

\begin{table}[h]
\centering
\caption{t5-small sem RGAT vs.\ t5-base com RGAT (fine-tuning completo).}
\label{tab:t5small-vs-rgat}
\begin{tabular}{lrrrl}
\toprule
\textbf{Configuração} & \textbf{exact match} & \textbf{exec} & \textbf{eval\_loss} & \textbf{tempo de treino} \\
\midrule
t5-base + RGAT (com \texttt{oncomx}, gc) & 0,0\% & 0,34\% & 1,6765 & 9h42min \\
t5-small \textbf{sem} RGAT (3 épocas)     & \textbf{1,38\%} & \textbf{1,37\%} & \textbf{1,1825} & \textbf{2h53min} \\
\bottomrule
\end{tabular}
\end{table}

O resultado é surpreendente: o modelo \emph{menor} e \emph{sem} qualquer
estrutura de grafo superou o modelo com RGAT em todas as métricas -- exact
match 4$\times$ maior (1,38\% vs.\ 0,34\% de exec no RGAT, comparação
direta de exec: 1,37\% vs.\ 0,34\%, também $\sim$4$\times$), \texttt{eval\_loss}
menor, e treino $3{,}4\times$ mais rápido (sem o custo de recomputação do
\emph{gradient checkpointing} nem a sobrecarga de propagar mensagens em
grafos de até 1300 nós por exemplo). É importante notar que essa
comparação \textbf{não é perfeitamente controlada}: t5-small e t5-base+RGAT
diferem tanto na presença do RGAT quanto no tamanho do modelo base, então
o experimento isola "RGAT vs.\ sem RGAT" apenas de forma aproximada, não
com uma variável por vez. Ainda assim, o resultado é informativo na direção
oposta à esperada -- se o RGAT estivesse agregando valor líquido, seria
razoável esperar que o modelo \emph{maior e com estrutura adicional}
superasse o menor e mais simples, não o contrário. Hipóteses para essa
inversão:

\begin{itemize}
\item O RGAT deste projeto foi desenhado e ajustado para os esquemas
pequenos e homogêneos do Spider ($\sim$5,3 tabelas/banco); grafos de
1000+ nós do ScienceBenchmark estão 2--3 ordens de grandeza acima da escala
para a qual a arquitetura foi validada, e a propagação de mensagens em
grafos tão densos pode diluir sinal útil em vez de reforçá-lo.
\item Um modelo maior e mais complexo (t5-base+RGAT) tipicamente precisa de
mais dados e/ou mais épocas para superar um modelo menor em fine-tuning --
com apenas 4732 exemplos e 3 épocas, o orçamento de treino pode ser
insuficiente para que a capacidade extra do RGAT se traduza em vantagem
prática, enquanto o t5-small, por ser mais simples, converge mais
rapidamente para um padrão razoável com o mesmo orçamento.
\item A complexidade adicional (mais parâmetros, mais uma fonte de
gradiente via o \emph{encoder} do grafo) pode estar introduzindo ruído de
otimização que atrapalha a convergência dentro do orçamento de treino
disponível, sem chegar a comprometer a estabilidade geral (o \emph{loss}
de treino do RGAT+t5-base ainda caiu de forma saudável, só não tão baixo
quanto seria necessário para competir com o t5-small nas métricas
discretas).
\end{itemize}

Este resultado não permite concluir que RGAT seja inerentemente inadequado
para text-to-SQL sobre esquemas grandes -- apenas que, \textbf{nas
condições específicas deste trabalho} (hardware de 8\,GB de VRAM, poucas
horas de treino disponíveis, 4732 exemplos), a complexidade adicional do
RGAT não se pagou em métricas discretas frente a um baseline T5 puro e
menor, treinado com o mesmo orçamento de tempo.

\subsection{Limitações metodológicas}
\label{sec:limitations}

Os dois experimentos de fine-tuning expuseram uma série de limitações,
tanto de hardware quanto de configuração, documentadas aqui por
transparência e como recomendação para trabalho futuro:

\begin{enumerate}
\item[(i)] \textbf{Ausência de split de teste genuíno.} O ScienceBenchmark,
como distribuído, só define \texttt{train} (seed+synth) e \texttt{dev}. O
mesmo split de \texttt{dev} usado para validação durante o treino
(seleção do "melhor" checkpoint) foi reutilizado como avaliação final --
não há um terceiro split, cego durante todo o processo, para reportar como
número de teste independente. O Spider, usado no treino-base, tem a mesma
limitação estrutural (sem test set público), então isso é consistente com
a metodologia herdada, mas permanece uma ressalva válida: o número final
reportado pode estar levemente inflado em relação a um teste verdadeiramente
cego.

\item[(ii)] \textbf{\texttt{torch.cuda.empty\_cache()} incondicional.} Uma
chamada incondicional a \texttt{torch.cuda.empty\_cache()} em
\texttt{seq2seq/models/graphix/rgat\_tuning.py::RGAT\_Layer.forward}, uma
vez por camada RGAT (12$\times$ por exemplo no t5-base), forçava uma
sincronização CUDA completa a cada chamada. Esse código é compartilhado
entre treino e avaliação (\texttt{modeling\_t5.py} importa a mesma classe
para ambos), então os resultados de zero-shot/PICARD/beam4 relatados acima
provavelmente rodaram mais devagar do que precisavam -- os números
continuam válidos, só o tempo de execução foi maior que o necessário. O
fix (só chamar \texttt{empty\_cache()} sob OOM real, não preventivamente)
foi validado por comparação direta: antes, nenhum passo de otimização
completava em mais de uma hora; depois, o mesmo passo completava em
minutos.

\item[(iii)] \textbf{\texttt{oncomx} tem esquema uniformemente grande
demais para 8\,GB de VRAM.} Todos os 1165 exemplos de treino de
\texttt{oncomx} produzem grafos RGAT entre 1022 e 1308 nós -- não existe
um subconjunto pequeno e seguro para amostrar. Mesmo os 20 exemplos
"menores" usados no few-shot (1022--1025 nós) ainda estão acima do teto de
900--1000 nós que se mostrou confiável durante o treino completo, e
provocaram OOMs individuais recorrentes durante o few-shot (capturados e
pulados por um mecanismo de segurança, não por acaso -- ver item v). Esta é
uma restrição de hardware, não uma limitação de amostragem: nenhuma
estratégia de seleção de exemplos resolve o problema sem reduzir
\texttt{max\_source\_length} (o que exigiria reprocessar toda a base de
grafos RGAT, fora do escopo deste trabalho). \textbf{Atualização}: esta
restrição foi contornada numa iteração posterior deste trabalho habilitando
corretamente o \emph{gradient checkpointing} do PyTorch, o que permitiu
incluir \texttt{oncomx} por completo no treino sem reduzir
\texttt{max\_source\_length} -- ver Seção~\ref{sec:gc-fix}.

\item[(iv)] \textbf{Ausência de \emph{warmup} de taxa de aprendizado no
fine-tuning completo, e espaçamento de checkpoint excessivo no few-shot.}
O fine-tuning completo usou \texttt{warmup\_ratio=0}, fazendo a taxa de
aprendizado saltar diretamente para $5\times10^{-5}$ no primeiro passo --
isso corresponde a um colapso temporário de geração (repetição de
n-gramas, ex.\ \texttt{"...loadversion = 301 and loadversion = 301 and..."}
até o limite de tokens), causando \texttt{eval\_exact\_match=0} em todos os
checkpoints intermediários mesmo com o \emph{loss} caindo normalmente. Um
\texttt{no\_repeat\_ngram\_size=3} foi adicionado durante o treino
(reduzindo o tempo de geração em $5\times$ e eliminando os loops de
n-grama), e o few-shot já nasceu com \texttt{warmup\_ratio=0,1} e
\texttt{no\_repeat\_ngram\_size=3} desde o início -- mas um padrão de
repetição mais sutil (alternância entre tokens SentencePiece de tamanhos
diferentes que decodificam para o mesmo caractere repetido, ex.\
\texttt{"6666666..."}) persistiu, sugerindo que o modelo continuou
sub-treinado para geração de literais numéricos mesmo ao final do processo.
Separadamente, no few-shot, o intervalo entre checkpoints
(\texttt{eval\_steps=47}) foi ajustado no meio do treino para reduzir o
tempo gasto em avaliação (cada rodada de eval levava até 111 minutos), mas
isso fez com que só um checkpoint intermediário (passo 47 de 90, 52\% do
treino) fosse salvo antes do fim -- o \texttt{load\_best\_model\_at\_end}
carregou esse checkpoint como "melhor" por ser o único candidato com
métrica registrada, e os pesos do passo 90 (o treino de fato mais avançado)
nunca chegaram a ser persistidos em disco.

\item[(v)] \textbf{Seleção de "melhor checkpoint" quebrada quando a métrica
empata.} O fine-tuning completo usou
\texttt{metric\_for\_best\_model="exact\_match"}. Como o colapso de geração
descrito em (iv) manteve \texttt{exact\_match=0,0} em \emph{todos} os
checkpoints intermediários, e a lógica de comparação do HuggingFace
Trainer só substitui o "melhor" registrado quando a nova métrica é
\emph{estritamente} maior, o primeiro checkpoint avaliado (passo 50 de
330, a versão \emph{menos} treinada de todas) nunca foi superado e acabou
sendo o modelo efetivamente salvo ao final do treino -- os pesos do passo
300 (97\% do treino) precisaram ser recuperados manualmente do diretório de
checkpoint correspondente para a avaliação real reportada na
Tabela~\ref{tab:finetuning-full}. O few-shot corrigiu isso usando
\texttt{metric\_for\_best\_model="eval\_loss"} com
\texttt{greater\_is\_better=false} (uma métrica contínua, que decresceu de
forma monotônica ao longo do treino), mas essa correção por si só não
evitou a limitação (iv) relacionada ao espaçamento de checkpoints.

\item[(vi)] \textbf{OOMs genuínos e travamento de driver de GPU.} Dois
crashes de OOM (\emph{out of memory}) genuínos interromperam o fine-tuning
completo no meio do treino, e um deles coincidiu com um travamento
completo do driver NVIDIA (\texttt{nvidia-smi} parou de reconhecer a
placa), exigindo reinicialização completa da máquina -- um reset de driver
isolado não foi suficiente para restaurar o desempenho normal da GPU.
Mitigações aplicadas: \emph{checkpointing} a cada 50 passos (em vez de 500,
o intervalo original herdado do treino no Spider, que nunca dispararia
dentro de um treino de 330 passos) e uma captura de OOM por micro-batch em
\texttt{seq2seq/utils/trainer.py::Seq2SeqTrainer.training\_step} que pula
apenas o exemplo problemático (zerando o gradiente parcialmente acumulado
e liberando cache) em vez de derrubar o processo inteiro -- validado
diretamente durante o few-shot, que sobreviveu a dezenas de OOMs
individuais de exemplos grandes de \texttt{oncomx} sem interrupção.
\end{enumerate}

\section{Comparação final}
\label{sec:final-comparison}

\begin{table}[h]
\centering
\caption{Comparação de todas as configurações avaliadas no dev set do ScienceBenchmark.}
\label{tab:final-comparison}
\begin{tabular}{lrrl}
\toprule
\textbf{Configuração} & \textbf{exact match} & \textbf{exec} & \textbf{Observação} \\
\midrule
Zero-shot (greedy, num\_beams=1)  & 9,0\%  & 8,87\% & referência \\
Zero-shot + beam search (n=4)      & \textbf{9,34\%} & \textbf{9,56\%} & melhor resultado geral \\
Zero-shot + PICARD                 & 9,0\%  & 8,87\% & PICARD inerte (falha de integração, \S\ref{sec:picard-attempt}) \\
Fine-tuning completo, sem \texttt{oncomx} (3 épocas)    & 0,0\%  & 0,34\% & pior que zero-shot (\S\ref{sec:finetuning-full}) \\
Fine-tuning few-shot (300 ex., 10 épocas) & 0,0\% & 0,0\% & pior que zero-shot (\S\ref{sec:finetuning-fewshot}) \\
Fine-tuning completo, com \texttt{oncomx} + gc (3 épocas) & 0,0\% & 0,34\% & eval\_loss $2,5\times$ menor, métricas discretas inalteradas (\S\ref{sec:gc-fix}) \\
Baseline t5-small sem RGAT (3 épocas)     & 1,38\% & 1,37\% & supera todos os fine-tunings com RGAT (\S\ref{sec:t5small-baseline}) \\
\bottomrule
\end{tabular}
\end{table}

O achado central deste conjunto de experimentos permanece, em essência,
\textbf{negativo em relação ao fine-tuning com RGAT}: em nenhuma das quatro
tentativas de ajuste fino (completo sem \texttt{oncomx}, \emph{few-shot},
completo com \texttt{oncomx} via \emph{gradient checkpointing}) o modelo
t5-base+RGAT superou o desempenho do checkpoint zero-shot com busca em
feixe nas métricas discretas -- mesmo depois de resolver a limitação de
VRAM que impedia incluir todos os domínios de treino. A correção do
\emph{gradient checkpointing} (\S\ref{sec:gc-fix}) trouxe uma melhora real
e mensurável na perda de avaliação (2,5$\times$ menor), mas essa melhora
não se propagou às métricas discretas -- um resultado que, por si só, já é
instrutivo: perda agregada e acerto exato/execução podem divergir
significativamente quando os erros do modelo são predominantemente
categóricos (nome de tabela/coluna errado) em vez de graduais.

O achado mais forte deste trabalho, no entanto, veio do baseline de
ablação (\S\ref{sec:t5small-baseline}): um \texttt{t5-small} \textbf{sem
nenhuma estrutura de grafo}, treinado com o mesmo orçamento de dados e
épocas, superou \emph{todas} as configurações com RGAT fine-tunadas neste
trabalho em exact match e exec, e o fez em menos de um terço do tempo de
treino. Combinado com o resultado do zero-shot (onde o RGAT+t5-base é a
única opção avaliada, pois não há um t5-small equivalente pré-treinado no
Spider para comparação direta), a evidência deste trabalho aponta para uma
conclusão qualificada: \textbf{nas condições de hardware e dados
disponíveis, a complexidade adicional do RGAT não demonstrou vantagem
sobre um T5 padrão no fine-tuning para o ScienceBenchmark} -- um resultado
negativo genuíno para a hipótese de que a estrutura de grafo ajudaria a
generalizar para esquemas maiores e de domínio distinto do treino
original. As limitações documentadas na Seção~\ref{sec:limitations}
(volume de dados, orçamento de épocas, ausência de um terceiro split de
teste cego) continuam válidas e sugerem que essa conclusão é específica
deste orçamento experimental, não necessariamente uma propriedade
fundamental da arquitetura RGAT -- mas dentro do escopo efetivamente
testado, o T5 puro foi tanto mais eficaz quanto mais barato de treinar.

% ou compilado standalone.

\section{Avaliação no ScienceBenchmark (zero-shot)}
\label{sec:sciencebenchmark-eval}

O checkpoint treinado no Spider (t5-base + RGAT, 5 épocas, 8577 exemplos de
treino, resultado no dev set do Spider: 41,9\% \emph{exact match} / 44,8\%
\emph{execution accuracy}) foi avaliado, sem nenhum re-treinamento, sobre os
299 exemplos do split de desenvolvimento do ScienceBenchmark
\cite{sciencebenchmark2023}, distribuídos em três bancos de domínios distintos:
\texttt{cordis} (projetos de pesquisa europeus), \texttt{oncomx} (biomedicina)
e \texttt{sdss} (astronomia).

\subsection{Resultado agregado}

\begin{table}[h]
\centering
\caption{Resultado do eval zero-shot do checkpoint Spider-treinado no ScienceBenchmark, por banco.}
\label{tab:sciencebenchmark-zeroshot}
\begin{tabular}{lrrrl}
\toprule
\textbf{Banco} & \textbf{n} & \textbf{exact match} & \textbf{exec} & \textbf{Exemplos avaliados} \\
\midrule
cordis            & 100 & 12,0\% & 15,0\% & 100/100 \\
oncomx            & 99  & 14,1\% & 11,1\% & 99/99 (exec: 97/99) \\
sdss              & 94  & 0,0\%  & 0,0\%  & 90/94\textsuperscript{a} \\
\midrule
\textbf{Total (agregado)} & \textbf{293} & \textbf{9,0\%} & \textbf{8,9\%} & exact: 289/293, exec: 287/293 \\
\bottomrule
\end{tabular}

\vspace{0.3em}
{\footnotesize\textsuperscript{a}4 exemplos do sdss descartados da métrica: parser
SQL do Spider (\texttt{third\_party/spider/process\_sql.py}) não suporta
parênteses agrupando condições no \texttt{WHERE}, um padrão ausente no SQL
gold do Spider mas presente no do ScienceBenchmark.}
\end{table}

Para efeito de comparação, o mesmo checkpoint avaliado no dev set do Spider
(1034 exemplos, mesmo domínio de treino) obteve 41,9\% exact match e 44,8\%
exec -- uma queda de $\sim$33 pontos percentuais ao mudar de domínio sem
re-treinamento, consistente com a expectativa de que arquiteturas
graph-aware treinadas em um conjunto de bancos pequenos e homogêneos (Spider,
$\sim$5,3 tabelas/banco em média) generalizam de forma limitada para bancos
maiores e de domínio mais específico (ScienceBenchmark, 19--25 tabelas no
\texttt{cordis}, colunas de medição numérica pouco usuais no \texttt{sdss}).

\subsection{Quebra por dificuldade}

Usando a classificação de dificuldade do avaliador oficial do Spider
(\emph{easy}/\emph{medium}/\emph{hard}/\emph{extra}, baseada em número de
componentes SQL, condições e agregações), o exact match cai a zero em
\emph{hard} e \emph{extra} nos três bancos, sem exceção:

\begin{table}[h]
\centering
\caption{Exact match por nível de dificuldade e banco.}
\label{tab:sciencebenchmark-hardness}
\begin{tabular}{lrrrr}
\toprule
\textbf{Banco} & \textbf{easy} & \textbf{medium} & \textbf{hard} & \textbf{extra} \\
\midrule
cordis (n=100) & 32,0\% (n=25) & 11,4\% (n=35) & 0\% (n=19) & 0\% (n=21) \\
oncomx (n=99)  & 42,9\% (n=21) & 13,9\% (n=36) & 0\% (n=22) & 0\% (n=20) \\
sdss (n=90)    & 0\% (n=6)     & 0\% (n=28)    & 0\% (n=20) & 0\% (n=36) \\
\bottomrule
\end{tabular}
\end{table}

O padrão observado sugere duas causas de queda distintas e sobrepostas: (i)
uma limitação de \emph{capacidade}, comum aos três bancos -- o modelo t5-base
resolve consultas simples mas não lida com consultas complexas (múltiplos
JOINs, subconsultas, agregações compostas) fora do domínio de treino; e (ii)
uma limitação de \emph{transferência de domínio}, específica do
\texttt{sdss} -- o resultado é zero mesmo nas consultas \emph{easy}, o que não
ocorre em \texttt{cordis}/\texttt{oncomx}, indicando que o vocabulário e a
estrutura de esquema do \texttt{sdss} (medições fotométricas, nomenclatura
astronômica) são suficientemente distantes do domínio de treino a ponto de
comprometer até os casos mais simples.

Uma observação metodológica reforça essa segunda leitura: o tamanho do
esquema (contagem de nós do grafo RGAT) e o banco de dados são quase
colineares nesta amostra -- \texttt{cordis}/\texttt{oncomx} sempre excedem
512 nós, \texttt{sdss} está sempre abaixo. Como \texttt{sdss} é o banco com o
\emph{menor} esquema em número de nós e ainda assim o único com resultado
nulo em todos os níveis de dificuldade, o fator determinante da queda não
parece ser o tamanho do esquema, e sim a distância de domínio/vocabulário em
relação ao Spider.

\subsection{Análise qualitativa dos erros}

Das 263 previsões incorretas (excluindo os 4 exemplos de \texttt{sdss} sem
gold parseável), 179 (68\%) foram inicialmente classificadas como ``SQL
sintaticamente inválido'' pelo parser oficial do Spider
(\texttt{third\_party/spider/process\_sql.py}) -- ou seja, o modelo nem
produziu uma consulta que o parser conseguisse interpretar como AST SQL.

\begin{table}[h]
\centering
\caption{Distribuição das categorias de erro (previsões incorretas, exact match = 0).}
\label{tab:error-categories}
\begin{tabular}{lrrrr}
\toprule
\textbf{Categoria} & \textbf{Total} & \textbf{cordis} & \textbf{oncomx} & \textbf{sdss} \\
\midrule
SQL sintaticamente inválido & 179 & 56 & 58 & 65 \\
Tabela errada                & 50  & 13 & 12 & 25 \\
Coluna errada                 & 19  & 9  & 10 & 0  \\
Outro erro semântico          & 10  & 5  & 5  & 0  \\
Agregação errada              & 5   & 5  & 0  & 0  \\
\bottomrule
\end{tabular}
\end{table}

Como o parser oficial do Spider já se mostrou limitado a uma gramática SQL
restrita neste trabalho (cf. o caso de parênteses agrupando condições no
\texttt{WHERE}, seção anterior), investigou-se se parte dessas 179 previsões
``inválidas'' seriam, na verdade, SQLite sintaticamente válido que o parser
simplesmente não reconhece. Cada uma das 179 previsões foi executada
diretamente contra o banco SQLite real correspondente (\texttt{sqlite3},
sem passar pelo parser AST do Spider). O resultado foi conclusivo:
\textbf{nenhuma das 179 (0\%) executou com sucesso} -- todas produziram um
erro real do SQLite, predominantemente \texttt{no such table} ou
\texttt{no such column} (por exemplo, \texttt{institutions.uncs\_id} em vez
de \texttt{unics\_id}, ou um join com a tabela inexistente \texttt{subjects}
em vez de \texttt{subject\_areas}). Isso descarta a hipótese de artefato de
parser para esta categoria: a maioria dos erros do modelo no
ScienceBenchmark é, de fato, alucinação de nomes de tabela/coluna
inexistentes no esquema ou SQL genuinamente malformado, não uma limitação da
ferramenta de avaliação.

\subsection{Decodificação restrita por gramática (PICARD)}
\label{sec:picard-attempt}

Dado que a maioria dos erros é SQL sintaticamente inválido, testou-se se a
decodificação restrita por gramática do PICARD \cite{picard2021}
-- que impede o modelo de gerar tokens que levariam a uma consulta
inválida -- reduziria essa taxa. O eval foi refeito com
\texttt{use\_picard=true}, \texttt{launch\_picard=true}, mantendo os demais
parâmetros idênticos.

\textbf{Resultado}: as 293 previsões geradas com PICARD ativo são
\emph{idênticas, token a token}, às geradas sem PICARD -- exact match e exec
saíram exatamente iguais aos valores da Tabela~\ref{tab:sciencebenchmark-zeroshot}.
Investigando a causa, constatou-se que o PICARD nunca chega a interceptar a
geração nesta configuração: em \texttt{seq2seq/models/graphix/rgat\_picard.py},
um ajuste anterior (necessário para corrigir a resolução da classe do modelo,
que sem ele carregava um T5 genérico em vez da variante com RGAT) passou a
importar a classe do modelo diretamente, contornando o mecanismo
(\texttt{model\_cls\_wrapper}) pelo qual o \texttt{generate()} restrito do
PICARD normalmente seria aplicado. O servidor PICARD chega a subir e ficar
ativo durante o eval, mas nunca é de fato consultado pela geração --
confirmado tanto pela ausência total de qualquer menção a
``picard''/``thrift'' nos logs do eval quanto pela identidade exata das
previsões. Trata-se de uma lacuna de integração pré-existente, não avaliada
anteriormente porque o eval com PICARD nunca havia sido exercitado de ponta
a ponta neste projeto; o conserto exigiria combinar os dois mecanismos de
substituição de classe (RGAT-aware e PICARD-aware) em vez de tratá-los como
alternativas mutuamente exclusivas, o que fica registrado aqui como limitação
conhecida e trabalho futuro.

\subsection{Busca em feixe (\emph{beam search}, num\_beams=4)}
\label{sec:beam4}

Como variação de decodificação adicional (mantendo o checkpoint zero-shot
inalterado), o eval foi refeito com \texttt{num\_beams=4} no lugar da
decodificação gulosa (\texttt{num\_beams=1}) usada nos resultados acima.

\textbf{Resultado}: 9,34\% exact match / 9,56\% exec (289 exemplos
pontuados para exact match, 287 para exec), \texttt{eval\_loss=0,6885} --
uma melhora pequena mas real sobre a decodificação gulosa (9,0\%/8,87\%),
consistente com o ganho usual, porém modesto, de busca em feixe sobre
geração gulosa em tarefas de geração estruturada.

\section{Fine-tuning no ScienceBenchmark}
\label{sec:finetuning}

A partir do mesmo checkpoint zero-shot, foram conduzidos dois experimentos de
ajuste fino sobre os dados de treino do ScienceBenchmark (4732 exemplos,
seed+synth combinados): um fine-tuning completo (todas as épocas sobre o
maior subconjunto de dados que a GPU disponível comporta) e um fine-tuning
\emph{few-shot} (poucas centenas de exemplos, mais épocas). Ambos foram
avaliados no mesmo split de dev usado no zero-shot, seguindo a mesma
metodologia de scoring (exact match + execução, mesmos filtros de
segurança contra travamentos de execução no \texttt{sdss}).

\textbf{Restrição de hardware determinante}: os experimentos rodaram numa
única GPU de consumidor (NVIDIA GeForce GTX 1070, 8\,GB de VRAM, sem tensor
cores), o mesmo hardware usado para os evals acima. Diferente do Spider
(esquemas pequenos, $\sim$5,3 tabelas/banco em média), os grafos RGAT do
ScienceBenchmark chegam a mais de 1300 nós por exemplo, o que por si só é
suficiente para saturar 8\,GB de VRAM durante o treino (fase que aloca
ativações para retropropagação, ao contrário do eval, que só faz
inferência).

\subsection{Fine-tuning completo}
\label{sec:finetuning-full}

\textbf{Configuração}: 3 épocas, \texttt{batch\_size=1} com
\texttt{gradient\_accumulation\_steps=32} (lote efetivo 32), taxa de
aprendizado $5\times10^{-5}$ com \texttt{adafactor}, 330 passos de
otimização totais. Por restrição de VRAM, exemplos cujo grafo RGAT excede
1000 nós foram descartados do treino -- isso excluiu \textbf{1165 dos 1165
exemplos de \texttt{oncomx}} presentes no split de treino após o primeiro
filtro (praticamente o banco inteiro; ver limitação (iii) abaixo), mantendo
\texttt{cordis} e \texttt{sdss} quase integralmente (3540 dos 4732 exemplos
totais permaneceram).

\begin{table}[h]
\centering
\caption{Fine-tuning completo: resultado final no dev set completo (289--293 exemplos, sem o filtro de tamanho de grafo aplicado ao treino).}
\label{tab:finetuning-full}
\begin{tabular}{lrrr}
\toprule
\textbf{Configuração} & \textbf{exact match} & \textbf{exec} & \textbf{eval\_loss} \\
\midrule
Zero-shot (referência) & 9,0\%  & 8,87\% & -- \\
Fine-tuning completo    & \textbf{0,0\%} & \textbf{0,34\%} & 4,2321 \\
\bottomrule
\end{tabular}
\end{table}

O resultado é \textbf{pior que o zero-shot}, apesar do \emph{loss} de
treino ter caído de forma saudável e monotônica até 0,469 ao final das 3
épocas. A causa raiz identificada: o modelo nunca viu nenhum exemplo de
\texttt{oncomx} durante o treino (excluído por VRAM), mas o dev set de
avaliação final inclui \texttt{oncomx} normalmente -- o desempenho ruim
nesse domínio nunca visto arrasta a média agregada para baixo, revertendo
o ganho que o fine-tuning trouxe em \texttt{cordis}/\texttt{sdss}. Esse é
um achado metodológico central deste trabalho: \textbf{fine-tuning
completo sob restrição severa de VRAM que exclui um dos domínios-alvo pode
piorar a generalização agregada}, mesmo quando todos os indicadores de
treino (perda) sugerem que o ajuste está funcionando.

\subsection{Fine-tuning \emph{few-shot}}
\label{sec:finetuning-fewshot}

Para tentar obter alguma representação de \texttt{oncomx} sem incorrer no
mesmo custo proibitivo de VRAM ao longo de um treino completo, foi montado
um segundo experimento com uma amostra pequena e deliberadamente balanceada
entre os três bancos: 140 exemplos de \texttt{cordis} + 140 de
\texttt{sdss} (amostragem aleatória) + 20 de \texttt{oncomx} (os 20 exemplos
de menor grafo disponíveis, 1022--1025 nós -- ver limitação (iii)), 300
exemplos no total, por 10 épocas (90 passos de otimização).

\begin{table}[h]
\centering
\caption{Fine-tuning few-shot: resultado final no dev set completo.}
\label{tab:finetuning-fewshot}
\begin{tabular}{lrrr}
\toprule
\textbf{Configuração} & \textbf{exact match} & \textbf{exec} & \textbf{eval\_loss} \\
\midrule
Zero-shot (referência)  & 9,0\%  & 8,87\% & -- \\
Fine-tuning completo     & 0,0\%  & 0,34\% & 4,2321 \\
Fine-tuning few-shot     & \textbf{0,0\%} & \textbf{0,0\%} & 4,8754 \\
\bottomrule
\end{tabular}
\end{table}

O resultado também ficou abaixo do zero-shot, e abaixo até do fine-tuning
completo no exec. Diferente do fine-tuning completo, aqui não há exclusão
sistemática de um banco (os três estão representados, ainda que
\texttt{oncomx} de forma mínima) -- a causa mais provável é volume de
dados insuficiente (300 exemplos) combinado com uma limitação metodológica
própria deste experimento, descrita em (iv) abaixo: o checkpoint final
salvo corresponde a apenas 52\% do treino planejado, não ao ponto de
convergência real.

\subsection{Correção do \emph{gradient checkpointing} e fine-tuning completo revisado}
\label{sec:gc-fix}

A limitação (iii) acima -- exclusão total de \texttt{oncomx} do treino por
falta de VRAM -- motivou uma segunda iteração deste trabalho, buscando
incluir os três domínios sem reduzir \texttt{max\_source\_length}. O
mecanismo natural para isso é o \emph{gradient checkpointing} do PyTorch
(\texttt{torch.utils.checkpoint}): em vez de manter em memória todas as
ativações intermediárias da rede para a retropropagação, ele as recalcula
sob demanda durante o \emph{backward pass}, trocando tempo de computação por
memória. A flag já existia em \texttt{configs/train.json}, mas estava sem
efeito para o modelo com RGAT -- dois problemas distintos, ambos
corrigidos:

\begin{itemize}
\item A classe \texttt{Model} em \texttt{seq2seq/models/graphix/rgat.py} é
um invólucro (\emph{wrapper}) que guarda o T5 com RGAT de fato como um
atributo interno (\texttt{self.pretrain\_model}), em vez de herdar sua
hierarquia de camadas. O método padrão
\texttt{gradient\_checkpointing\_enable()} do HuggingFace percorre
recursivamente as submódulos de uma classe para ativar o checkpointing em
cada um -- como o RGAT vive dentro de um atributo e não na hierarquia
percorrida, a ativação nunca chegava lá. Corrigido delegando
explicitamente a chamada ao \texttt{pretrain\_model} interno.
\item Mesmo com a delegação corrigida, o \texttt{T5Stack.forward()} (em
\texttt{seq2seq/models/modeling\_t5.py}) constrói, quando o checkpointing
está ativo, uma função de encapsulamento (\emph{closure}) para cada bloco
que é passada a \texttt{torch.utils.checkpoint.checkpoint()} -- e essa
closure não repassava \texttt{graph\_batch}/\texttt{relation\_emb} ao
bloco, fazendo o RGAT ser silenciosamente ignorado (ou falhar) sempre que
o checkpointing estava ativo. Corrigido capturando esses dois objetos por
\emph{closure} do Python (não são tensores diferenciáveis que o mecanismo
de checkpoint precise gerenciar, então isso é seguro tanto na chamada
inicial quanto no recálculo do \emph{backward}).
\end{itemize}

A correção foi validada com um teste de fumaça isolado sobre os seis
maiores grafos de \texttt{oncomx} (1217--1308 nós, os que mais
frequentemente causavam OOM antes), confirmando treino sem erros e perda
decrescente normalmente.

Com o \emph{gradient checkpointing} funcionando, o fine-tuning completo foi
refeito com os mesmos hiperparâmetros da Seção~\ref{sec:finetuning-full}
(3 épocas, \texttt{batch\_size=1}, \texttt{gradient\_accumulation\_steps=32},
$5\times10^{-5}$ com \texttt{adafactor}), desta vez sem nenhuma exclusão de
exemplo por tamanho de grafo -- os 4732 exemplos de treino (incluindo
\texttt{oncomx} por completo) participaram do treino.

\begin{table}[h]
\centering
\caption{Fine-tuning completo, antes e depois da correção do gradient checkpointing.}
\label{tab:finetuning-full-gc}
\begin{tabular}{lrrrl}
\toprule
\textbf{Configuração} & \textbf{exact match} & \textbf{exec} & \textbf{eval\_loss} & \textbf{train\_loss / tempo} \\
\midrule
Fine-tuning completo (sem \texttt{oncomx}) & 0,0\% & 0,34\% & 4,2321 & 0,469 / -- \\
Fine-tuning completo (com \texttt{oncomx}, gc) & 0,0\% & \textbf{0,34\%} & \textbf{1,6765} & 2,215 / 9h42min \\
\bottomrule
\end{tabular}
\end{table}

O resultado é um achado genuinamente nuançado: incluir \texttt{oncomx}
reduziu o \texttt{eval\_loss} em quase $2{,}5\times$ (4,2321 $\to$ 1,6765) --
uma melhora clara e esperada, já que o modelo agora vê exemplos do domínio
que antes lhe era totalmente estranho. Mas essa melhora \textbf{não se
traduziu em nenhum ganho nas métricas discretas}: \texttt{exact\_match}
permanece em 0\% e \texttt{exec} praticamente inalterado (0,34\%). Isso
sugere que a hipótese da Seção~\ref{sec:finetuning-full} (exclusão de
\texttt{oncomx} como causa única do resultado ruim) era \emph{necessária
mas não suficiente}: mesmo com todos os três domínios representados no
treino, os erros do modelo continuam sendo do tipo que quebra
\texttt{exact\_match}/\texttt{exec} de forma binária -- nome de
tabela/coluna incorreto, estrutura de \texttt{JOIN} equivocada -- e uma
perda agregada mais baixa (que mede o quão "próxima" a distribuição de
tokens gerada está do alvo, token a token) não impede esse tipo de erro
categórico. As causas mais prováveis são um orçamento de treino ainda
pequeno (3 épocas sobre 4732 exemplos heterogêneos, contra 5 épocas sobre
8577 exemplos homogêneos no Spider original) e/ou uma limitação de
capacidade do próprio t5-base diante da diversidade de esquemas do
ScienceBenchmark, discutida em conjunto com o achado da próxima seção.

\subsection{Baseline t5-small sem RGAT}
\label{sec:t5small-baseline}

Para isolar o efeito específico do RGAT dos demais fatores (tamanho do
modelo, orçamento de treino, dados), foi conduzido um experimento de
ablação: um T5 padrão, \textbf{sem nenhuma camada RGAT ou estrutura de
grafo}, na variante \texttt{t5-small} (60M parâmetros, contra os $\sim$220M
do \texttt{t5-base} usado em todos os experimentos anteriores). O restante
do pipeline foi mantido idêntico -- mesmo pré-processamento (a mesma
serialização de esquema textual produzida pelas etapas de grafo, ainda que
o grafo em si não seja consumido), mesmo split de treino/dev (4732/299
exemplos), mesmos hiperparâmetros (3 épocas, \texttt{batch\_size=1},
\texttt{gradient\_accumulation\_steps=32}, $5\times10^{-5}$ com
\texttt{adafactor}). Uma nova classe \texttt{Model}
(\texttt{seq2seq/models/graphix/plain.py}) carrega um
\texttt{T5ForConditionalGeneration} padrão do HuggingFace e ignora
completamente o \texttt{graph\_batch}, selecionável por uma variável de
ambiente (\texttt{GRAPHIX\_MODEL\_VARIANT=plain}) sem duplicar o restante
do script de treino.

\begin{table}[h]
\centering
\caption{t5-small sem RGAT vs.\ t5-base com RGAT (fine-tuning completo).}
\label{tab:t5small-vs-rgat}
\begin{tabular}{lrrrl}
\toprule
\textbf{Configuração} & \textbf{exact match} & \textbf{exec} & \textbf{eval\_loss} & \textbf{tempo de treino} \\
\midrule
t5-base + RGAT (com \texttt{oncomx}, gc) & 0,0\% & 0,34\% & 1,6765 & 9h42min \\
t5-small \textbf{sem} RGAT (3 épocas)     & \textbf{1,38\%} & \textbf{1,37\%} & \textbf{1,1825} & \textbf{2h53min} \\
\bottomrule
\end{tabular}
\end{table}

O resultado é surpreendente: o modelo \emph{menor} e \emph{sem} qualquer
estrutura de grafo superou o modelo com RGAT em todas as métricas -- exact
match 4$\times$ maior (1,38\% vs.\ 0,34\% de exec no RGAT, comparação
direta de exec: 1,37\% vs.\ 0,34\%, também $\sim$4$\times$), \texttt{eval\_loss}
menor, e treino $3{,}4\times$ mais rápido (sem o custo de recomputação do
\emph{gradient checkpointing} nem a sobrecarga de propagar mensagens em
grafos de até 1300 nós por exemplo). É importante notar que essa
comparação \textbf{não é perfeitamente controlada}: t5-small e t5-base+RGAT
diferem tanto na presença do RGAT quanto no tamanho do modelo base, então
o experimento isola "RGAT vs.\ sem RGAT" apenas de forma aproximada, não
com uma variável por vez. Ainda assim, o resultado é informativo na direção
oposta à esperada -- se o RGAT estivesse agregando valor líquido, seria
razoável esperar que o modelo \emph{maior e com estrutura adicional}
superasse o menor e mais simples, não o contrário. Hipóteses para essa
inversão:

\begin{itemize}
\item O RGAT deste projeto foi desenhado e ajustado para os esquemas
pequenos e homogêneos do Spider ($\sim$5,3 tabelas/banco); grafos de
1000+ nós do ScienceBenchmark estão 2--3 ordens de grandeza acima da escala
para a qual a arquitetura foi validada, e a propagação de mensagens em
grafos tão densos pode diluir sinal útil em vez de reforçá-lo.
\item Um modelo maior e mais complexo (t5-base+RGAT) tipicamente precisa de
mais dados e/ou mais épocas para superar um modelo menor em fine-tuning --
com apenas 4732 exemplos e 3 épocas, o orçamento de treino pode ser
insuficiente para que a capacidade extra do RGAT se traduza em vantagem
prática, enquanto o t5-small, por ser mais simples, converge mais
rapidamente para um padrão razoável com o mesmo orçamento.
\item A complexidade adicional (mais parâmetros, mais uma fonte de
gradiente via o \emph{encoder} do grafo) pode estar introduzindo ruído de
otimização que atrapalha a convergência dentro do orçamento de treino
disponível, sem chegar a comprometer a estabilidade geral (o \emph{loss}
de treino do RGAT+t5-base ainda caiu de forma saudável, só não tão baixo
quanto seria necessário para competir com o t5-small nas métricas
discretas).
\end{itemize}

Este resultado não permite concluir que RGAT seja inerentemente inadequado
para text-to-SQL sobre esquemas grandes -- apenas que, \textbf{nas
condições específicas deste trabalho} (hardware de 8\,GB de VRAM, poucas
horas de treino disponíveis, 4732 exemplos), a complexidade adicional do
RGAT não se pagou em métricas discretas frente a um baseline T5 puro e
menor, treinado com o mesmo orçamento de tempo.

\subsection{Limitações metodológicas}
\label{sec:limitations}

Os dois experimentos de fine-tuning expuseram uma série de limitações,
tanto de hardware quanto de configuração, documentadas aqui por
transparência e como recomendação para trabalho futuro:

\begin{enumerate}
\item[(i)] \textbf{Ausência de split de teste genuíno.} O ScienceBenchmark,
como distribuído, só define \texttt{train} (seed+synth) e \texttt{dev}. O
mesmo split de \texttt{dev} usado para validação durante o treino
(seleção do "melhor" checkpoint) foi reutilizado como avaliação final --
não há um terceiro split, cego durante todo o processo, para reportar como
número de teste independente. O Spider, usado no treino-base, tem a mesma
limitação estrutural (sem test set público), então isso é consistente com
a metodologia herdada, mas permanece uma ressalva válida: o número final
reportado pode estar levemente inflado em relação a um teste verdadeiramente
cego.

\item[(ii)] \textbf{\texttt{torch.cuda.empty\_cache()} incondicional.} Uma
chamada incondicional a \texttt{torch.cuda.empty\_cache()} em
\texttt{seq2seq/models/graphix/rgat\_tuning.py::RGAT\_Layer.forward}, uma
vez por camada RGAT (12$\times$ por exemplo no t5-base), forçava uma
sincronização CUDA completa a cada chamada. Esse código é compartilhado
entre treino e avaliação (\texttt{modeling\_t5.py} importa a mesma classe
para ambos), então os resultados de zero-shot/PICARD/beam4 relatados acima
provavelmente rodaram mais devagar do que precisavam -- os números
continuam válidos, só o tempo de execução foi maior que o necessário. O
fix (só chamar \texttt{empty\_cache()} sob OOM real, não preventivamente)
foi validado por comparação direta: antes, nenhum passo de otimização
completava em mais de uma hora; depois, o mesmo passo completava em
minutos.

\item[(iii)] \textbf{\texttt{oncomx} tem esquema uniformemente grande
demais para 8\,GB de VRAM.} Todos os 1165 exemplos de treino de
\texttt{oncomx} produzem grafos RGAT entre 1022 e 1308 nós -- não existe
um subconjunto pequeno e seguro para amostrar. Mesmo os 20 exemplos
"menores" usados no few-shot (1022--1025 nós) ainda estão acima do teto de
900--1000 nós que se mostrou confiável durante o treino completo, e
provocaram OOMs individuais recorrentes durante o few-shot (capturados e
pulados por um mecanismo de segurança, não por acaso -- ver item v). Esta é
uma restrição de hardware, não uma limitação de amostragem: nenhuma
estratégia de seleção de exemplos resolve o problema sem reduzir
\texttt{max\_source\_length} (o que exigiria reprocessar toda a base de
grafos RGAT, fora do escopo deste trabalho). \textbf{Atualização}: esta
restrição foi contornada numa iteração posterior deste trabalho habilitando
corretamente o \emph{gradient checkpointing} do PyTorch, o que permitiu
incluir \texttt{oncomx} por completo no treino sem reduzir
\texttt{max\_source\_length} -- ver Seção~\ref{sec:gc-fix}.

\item[(iv)] \textbf{Ausência de \emph{warmup} de taxa de aprendizado no
fine-tuning completo, e espaçamento de checkpoint excessivo no few-shot.}
O fine-tuning completo usou \texttt{warmup\_ratio=0}, fazendo a taxa de
aprendizado saltar diretamente para $5\times10^{-5}$ no primeiro passo --
isso corresponde a um colapso temporário de geração (repetição de
n-gramas, ex.\ \texttt{"...loadversion = 301 and loadversion = 301 and..."}
até o limite de tokens), causando \texttt{eval\_exact\_match=0} em todos os
checkpoints intermediários mesmo com o \emph{loss} caindo normalmente. Um
\texttt{no\_repeat\_ngram\_size=3} foi adicionado durante o treino
(reduzindo o tempo de geração em $5\times$ e eliminando os loops de
n-grama), e o few-shot já nasceu com \texttt{warmup\_ratio=0,1} e
\texttt{no\_repeat\_ngram\_size=3} desde o início -- mas um padrão de
repetição mais sutil (alternância entre tokens SentencePiece de tamanhos
diferentes que decodificam para o mesmo caractere repetido, ex.\
\texttt{"6666666..."}) persistiu, sugerindo que o modelo continuou
sub-treinado para geração de literais numéricos mesmo ao final do processo.
Separadamente, no few-shot, o intervalo entre checkpoints
(\texttt{eval\_steps=47}) foi ajustado no meio do treino para reduzir o
tempo gasto em avaliação (cada rodada de eval levava até 111 minutos), mas
isso fez com que só um checkpoint intermediário (passo 47 de 90, 52\% do
treino) fosse salvo antes do fim -- o \texttt{load\_best\_model\_at\_end}
carregou esse checkpoint como "melhor" por ser o único candidato com
métrica registrada, e os pesos do passo 90 (o treino de fato mais avançado)
nunca chegaram a ser persistidos em disco.

\item[(v)] \textbf{Seleção de "melhor checkpoint" quebrada quando a métrica
empata.} O fine-tuning completo usou
\texttt{metric\_for\_best\_model="exact\_match"}. Como o colapso de geração
descrito em (iv) manteve \texttt{exact\_match=0,0} em \emph{todos} os
checkpoints intermediários, e a lógica de comparação do HuggingFace
Trainer só substitui o "melhor" registrado quando a nova métrica é
\emph{estritamente} maior, o primeiro checkpoint avaliado (passo 50 de
330, a versão \emph{menos} treinada de todas) nunca foi superado e acabou
sendo o modelo efetivamente salvo ao final do treino -- os pesos do passo
300 (97\% do treino) precisaram ser recuperados manualmente do diretório de
checkpoint correspondente para a avaliação real reportada na
Tabela~\ref{tab:finetuning-full}. O few-shot corrigiu isso usando
\texttt{metric\_for\_best\_model="eval\_loss"} com
\texttt{greater\_is\_better=false} (uma métrica contínua, que decresceu de
forma monotônica ao longo do treino), mas essa correção por si só não
evitou a limitação (iv) relacionada ao espaçamento de checkpoints.

\item[(vi)] \textbf{OOMs genuínos e travamento de driver de GPU.} Dois
crashes de OOM (\emph{out of memory}) genuínos interromperam o fine-tuning
completo no meio do treino, e um deles coincidiu com um travamento
completo do driver NVIDIA (\texttt{nvidia-smi} parou de reconhecer a
placa), exigindo reinicialização completa da máquina -- um reset de driver
isolado não foi suficiente para restaurar o desempenho normal da GPU.
Mitigações aplicadas: \emph{checkpointing} a cada 50 passos (em vez de 500,
o intervalo original herdado do treino no Spider, que nunca dispararia
dentro de um treino de 330 passos) e uma captura de OOM por micro-batch em
\texttt{seq2seq/utils/trainer.py::Seq2SeqTrainer.training\_step} que pula
apenas o exemplo problemático (zerando o gradiente parcialmente acumulado
e liberando cache) em vez de derrubar o processo inteiro -- validado
diretamente durante o few-shot, que sobreviveu a dezenas de OOMs
individuais de exemplos grandes de \texttt{oncomx} sem interrupção.
\end{enumerate}

\section{Comparação final}
\label{sec:final-comparison}

\begin{table}[h]
\centering
\caption{Comparação de todas as configurações avaliadas no dev set do ScienceBenchmark.}
\label{tab:final-comparison}
\begin{tabular}{lrrl}
\toprule
\textbf{Configuração} & \textbf{exact match} & \textbf{exec} & \textbf{Observação} \\
\midrule
Zero-shot (greedy, num\_beams=1)  & 9,0\%  & 8,87\% & referência \\
Zero-shot + beam search (n=4)      & \textbf{9,34\%} & \textbf{9,56\%} & melhor resultado geral \\
Zero-shot + PICARD                 & 9,0\%  & 8,87\% & PICARD inerte (falha de integração, \S\ref{sec:picard-attempt}) \\
Fine-tuning completo, sem \texttt{oncomx} (3 épocas)    & 0,0\%  & 0,34\% & pior que zero-shot (\S\ref{sec:finetuning-full}) \\
Fine-tuning few-shot (300 ex., 10 épocas) & 0,0\% & 0,0\% & pior que zero-shot (\S\ref{sec:finetuning-fewshot}) \\
Fine-tuning completo, com \texttt{oncomx} + gc (3 épocas) & 0,0\% & 0,34\% & eval\_loss $2,5\times$ menor, métricas discretas inalteradas (\S\ref{sec:gc-fix}) \\
Baseline t5-small sem RGAT (3 épocas)     & 1,38\% & 1,37\% & supera todos os fine-tunings com RGAT (\S\ref{sec:t5small-baseline}) \\
\bottomrule
\end{tabular}
\end{table}

O achado central deste conjunto de experimentos permanece, em essência,
\textbf{negativo em relação ao fine-tuning com RGAT}: em nenhuma das quatro
tentativas de ajuste fino (completo sem \texttt{oncomx}, \emph{few-shot},
completo com \texttt{oncomx} via \emph{gradient checkpointing}) o modelo
t5-base+RGAT superou o desempenho do checkpoint zero-shot com busca em
feixe nas métricas discretas -- mesmo depois de resolver a limitação de
VRAM que impedia incluir todos os domínios de treino. A correção do
\emph{gradient checkpointing} (\S\ref{sec:gc-fix}) trouxe uma melhora real
e mensurável na perda de avaliação (2,5$\times$ menor), mas essa melhora
não se propagou às métricas discretas -- um resultado que, por si só, já é
instrutivo: perda agregada e acerto exato/execução podem divergir
significativamente quando os erros do modelo são predominantemente
categóricos (nome de tabela/coluna errado) em vez de graduais.

O achado mais forte deste trabalho, no entanto, veio do baseline de
ablação (\S\ref{sec:t5small-baseline}): um \texttt{t5-small} \textbf{sem
nenhuma estrutura de grafo}, treinado com o mesmo orçamento de dados e
épocas, superou \emph{todas} as configurações com RGAT fine-tunadas neste
trabalho em exact match e exec, e o fez em menos de um terço do tempo de
treino. Combinado com o resultado do zero-shot (onde o RGAT+t5-base é a
única opção avaliada, pois não há um t5-small equivalente pré-treinado no
Spider para comparação direta), a evidência deste trabalho aponta para uma
conclusão qualificada: \textbf{nas condições de hardware e dados
disponíveis, a complexidade adicional do RGAT não demonstrou vantagem
sobre um T5 padrão no fine-tuning para o ScienceBenchmark} -- um resultado
negativo genuíno para a hipótese de que a estrutura de grafo ajudaria a
generalizar para esquemas maiores e de domínio distinto do treino
original. As limitações documentadas na Seção~\ref{sec:limitations}
(volume de dados, orçamento de épocas, ausência de um terceiro split de
teste cego) continuam válidas e sugerem que essa conclusão é específica
deste orçamento experimental, não necessariamente uma propriedade
fundamental da arquitetura RGAT -- mas dentro do escopo efetivamente
testado, o T5 puro foi tanto mais eficaz quanto mais barato de treinar.

% ou compilado standalone.

\section{Avaliação no ScienceBenchmark (zero-shot)}
\label{sec:sciencebenchmark-eval}

O checkpoint treinado no Spider (t5-base + RGAT, 5 épocas, 8577 exemplos de
treino, resultado no dev set do Spider: 41,9\% \emph{exact match} / 44,8\%
\emph{execution accuracy}) foi avaliado, sem nenhum re-treinamento, sobre os
299 exemplos do split de desenvolvimento do ScienceBenchmark
\cite{sciencebenchmark2023}, distribuídos em três bancos de domínios distintos:
\texttt{cordis} (projetos de pesquisa europeus), \texttt{oncomx} (biomedicina)
e \texttt{sdss} (astronomia).

\subsection{Resultado agregado}

\begin{table}[h]
\centering
\caption{Resultado do eval zero-shot do checkpoint Spider-treinado no ScienceBenchmark, por banco.}
\label{tab:sciencebenchmark-zeroshot}
\begin{tabular}{lrrrl}
\toprule
\textbf{Banco} & \textbf{n} & \textbf{exact match} & \textbf{exec} & \textbf{Exemplos avaliados} \\
\midrule
cordis            & 100 & 12,0\% & 15,0\% & 100/100 \\
oncomx            & 99  & 14,1\% & 11,1\% & 99/99 (exec: 97/99) \\
sdss              & 94  & 0,0\%  & 0,0\%  & 90/94\textsuperscript{a} \\
\midrule
\textbf{Total (agregado)} & \textbf{293} & \textbf{9,0\%} & \textbf{8,9\%} & exact: 289/293, exec: 287/293 \\
\bottomrule
\end{tabular}

\vspace{0.3em}
{\footnotesize\textsuperscript{a}4 exemplos do sdss descartados da métrica: parser
SQL do Spider (\texttt{third\_party/spider/process\_sql.py}) não suporta
parênteses agrupando condições no \texttt{WHERE}, um padrão ausente no SQL
gold do Spider mas presente no do ScienceBenchmark.}
\end{table}

Para efeito de comparação, o mesmo checkpoint avaliado no dev set do Spider
(1034 exemplos, mesmo domínio de treino) obteve 41,9\% exact match e 44,8\%
exec -- uma queda de $\sim$33 pontos percentuais ao mudar de domínio sem
re-treinamento, consistente com a expectativa de que arquiteturas
graph-aware treinadas em um conjunto de bancos pequenos e homogêneos (Spider,
$\sim$5,3 tabelas/banco em média) generalizam de forma limitada para bancos
maiores e de domínio mais específico (ScienceBenchmark, 19--25 tabelas no
\texttt{cordis}, colunas de medição numérica pouco usuais no \texttt{sdss}).

\subsection{Quebra por dificuldade}

Usando a classificação de dificuldade do avaliador oficial do Spider
(\emph{easy}/\emph{medium}/\emph{hard}/\emph{extra}, baseada em número de
componentes SQL, condições e agregações), o exact match cai a zero em
\emph{hard} e \emph{extra} nos três bancos, sem exceção:

\begin{table}[h]
\centering
\caption{Exact match por nível de dificuldade e banco.}
\label{tab:sciencebenchmark-hardness}
\begin{tabular}{lrrrr}
\toprule
\textbf{Banco} & \textbf{easy} & \textbf{medium} & \textbf{hard} & \textbf{extra} \\
\midrule
cordis (n=100) & 32,0\% (n=25) & 11,4\% (n=35) & 0\% (n=19) & 0\% (n=21) \\
oncomx (n=99)  & 42,9\% (n=21) & 13,9\% (n=36) & 0\% (n=22) & 0\% (n=20) \\
sdss (n=90)    & 0\% (n=6)     & 0\% (n=28)    & 0\% (n=20) & 0\% (n=36) \\
\bottomrule
\end{tabular}
\end{table}

O padrão observado sugere duas causas de queda distintas e sobrepostas: (i)
uma limitação de \emph{capacidade}, comum aos três bancos -- o modelo t5-base
resolve consultas simples mas não lida com consultas complexas (múltiplos
JOINs, subconsultas, agregações compostas) fora do domínio de treino; e (ii)
uma limitação de \emph{transferência de domínio}, específica do
\texttt{sdss} -- o resultado é zero mesmo nas consultas \emph{easy}, o que não
ocorre em \texttt{cordis}/\texttt{oncomx}, indicando que o vocabulário e a
estrutura de esquema do \texttt{sdss} (medições fotométricas, nomenclatura
astronômica) são suficientemente distantes do domínio de treino a ponto de
comprometer até os casos mais simples.

Uma observação metodológica reforça essa segunda leitura: o tamanho do
esquema (contagem de nós do grafo RGAT) e o banco de dados são quase
colineares nesta amostra -- \texttt{cordis}/\texttt{oncomx} sempre excedem
512 nós, \texttt{sdss} está sempre abaixo. Como \texttt{sdss} é o banco com o
\emph{menor} esquema em número de nós e ainda assim o único com resultado
nulo em todos os níveis de dificuldade, o fator determinante da queda não
parece ser o tamanho do esquema, e sim a distância de domínio/vocabulário em
relação ao Spider.

\subsection{Análise qualitativa dos erros}

Das 263 previsões incorretas (excluindo os 4 exemplos de \texttt{sdss} sem
gold parseável), 179 (68\%) foram inicialmente classificadas como ``SQL
sintaticamente inválido'' pelo parser oficial do Spider
(\texttt{third\_party/spider/process\_sql.py}) -- ou seja, o modelo nem
produziu uma consulta que o parser conseguisse interpretar como AST SQL.

\begin{table}[h]
\centering
\caption{Distribuição das categorias de erro (previsões incorretas, exact match = 0).}
\label{tab:error-categories}
\begin{tabular}{lrrrr}
\toprule
\textbf{Categoria} & \textbf{Total} & \textbf{cordis} & \textbf{oncomx} & \textbf{sdss} \\
\midrule
SQL sintaticamente inválido & 179 & 56 & 58 & 65 \\
Tabela errada                & 50  & 13 & 12 & 25 \\
Coluna errada                 & 19  & 9  & 10 & 0  \\
Outro erro semântico          & 10  & 5  & 5  & 0  \\
Agregação errada              & 5   & 5  & 0  & 0  \\
\bottomrule
\end{tabular}
\end{table}

Como o parser oficial do Spider já se mostrou limitado a uma gramática SQL
restrita neste trabalho (cf. o caso de parênteses agrupando condições no
\texttt{WHERE}, seção anterior), investigou-se se parte dessas 179 previsões
``inválidas'' seriam, na verdade, SQLite sintaticamente válido que o parser
simplesmente não reconhece. Cada uma das 179 previsões foi executada
diretamente contra o banco SQLite real correspondente (\texttt{sqlite3},
sem passar pelo parser AST do Spider). O resultado foi conclusivo:
\textbf{nenhuma das 179 (0\%) executou com sucesso} -- todas produziram um
erro real do SQLite, predominantemente \texttt{no such table} ou
\texttt{no such column} (por exemplo, \texttt{institutions.uncs\_id} em vez
de \texttt{unics\_id}, ou um join com a tabela inexistente \texttt{subjects}
em vez de \texttt{subject\_areas}). Isso descarta a hipótese de artefato de
parser para esta categoria: a maioria dos erros do modelo no
ScienceBenchmark é, de fato, alucinação de nomes de tabela/coluna
inexistentes no esquema ou SQL genuinamente malformado, não uma limitação da
ferramenta de avaliação.

\subsection{Decodificação restrita por gramática (PICARD)}
\label{sec:picard-attempt}

Dado que a maioria dos erros é SQL sintaticamente inválido, testou-se se a
decodificação restrita por gramática do PICARD \cite{picard2021}
-- que impede o modelo de gerar tokens que levariam a uma consulta
inválida -- reduziria essa taxa. O eval foi refeito com
\texttt{use\_picard=true}, \texttt{launch\_picard=true}, mantendo os demais
parâmetros idênticos.

\textbf{Resultado}: as 293 previsões geradas com PICARD ativo são
\emph{idênticas, token a token}, às geradas sem PICARD -- exact match e exec
saíram exatamente iguais aos valores da Tabela~\ref{tab:sciencebenchmark-zeroshot}.
Investigando a causa, constatou-se que o PICARD nunca chega a interceptar a
geração nesta configuração: em \texttt{seq2seq/models/graphix/rgat\_picard.py},
um ajuste anterior (necessário para corrigir a resolução da classe do modelo,
que sem ele carregava um T5 genérico em vez da variante com RGAT) passou a
importar a classe do modelo diretamente, contornando o mecanismo
(\texttt{model\_cls\_wrapper}) pelo qual o \texttt{generate()} restrito do
PICARD normalmente seria aplicado. O servidor PICARD chega a subir e ficar
ativo durante o eval, mas nunca é de fato consultado pela geração --
confirmado tanto pela ausência total de qualquer menção a
``picard''/``thrift'' nos logs do eval quanto pela identidade exata das
previsões. Trata-se de uma lacuna de integração pré-existente, não avaliada
anteriormente porque o eval com PICARD nunca havia sido exercitado de ponta
a ponta neste projeto; o conserto exigiria combinar os dois mecanismos de
substituição de classe (RGAT-aware e PICARD-aware) em vez de tratá-los como
alternativas mutuamente exclusivas, o que fica registrado aqui como limitação
conhecida e trabalho futuro.

\subsection{Busca em feixe (\emph{beam search}, num\_beams=4)}
\label{sec:beam4}

Como variação de decodificação adicional (mantendo o checkpoint zero-shot
inalterado), o eval foi refeito com \texttt{num\_beams=4} no lugar da
decodificação gulosa (\texttt{num\_beams=1}) usada nos resultados acima.

\textbf{Resultado}: 9,34\% exact match / 9,56\% exec (289 exemplos
pontuados para exact match, 287 para exec), \texttt{eval\_loss=0,6885} --
uma melhora pequena mas real sobre a decodificação gulosa (9,0\%/8,87\%),
consistente com o ganho usual, porém modesto, de busca em feixe sobre
geração gulosa em tarefas de geração estruturada.

\section{Fine-tuning no ScienceBenchmark}
\label{sec:finetuning}

A partir do mesmo checkpoint zero-shot, foram conduzidos dois experimentos de
ajuste fino sobre os dados de treino do ScienceBenchmark (4732 exemplos,
seed+synth combinados): um fine-tuning completo (todas as épocas sobre o
maior subconjunto de dados que a GPU disponível comporta) e um fine-tuning
\emph{few-shot} (poucas centenas de exemplos, mais épocas). Ambos foram
avaliados no mesmo split de dev usado no zero-shot, seguindo a mesma
metodologia de scoring (exact match + execução, mesmos filtros de
segurança contra travamentos de execução no \texttt{sdss}).

\textbf{Restrição de hardware determinante}: os experimentos rodaram numa
única GPU de consumidor (NVIDIA GeForce GTX 1070, 8\,GB de VRAM, sem tensor
cores), o mesmo hardware usado para os evals acima. Diferente do Spider
(esquemas pequenos, $\sim$5,3 tabelas/banco em média), os grafos RGAT do
ScienceBenchmark chegam a mais de 1300 nós por exemplo, o que por si só é
suficiente para saturar 8\,GB de VRAM durante o treino (fase que aloca
ativações para retropropagação, ao contrário do eval, que só faz
inferência).

\subsection{Fine-tuning completo}
\label{sec:finetuning-full}

\textbf{Configuração}: 3 épocas, \texttt{batch\_size=1} com
\texttt{gradient\_accumulation\_steps=32} (lote efetivo 32), taxa de
aprendizado $5\times10^{-5}$ com \texttt{adafactor}, 330 passos de
otimização totais. Por restrição de VRAM, exemplos cujo grafo RGAT excede
1000 nós foram descartados do treino -- isso excluiu \textbf{1165 dos 1165
exemplos de \texttt{oncomx}} presentes no split de treino após o primeiro
filtro (praticamente o banco inteiro; ver limitação (iii) abaixo), mantendo
\texttt{cordis} e \texttt{sdss} quase integralmente (3540 dos 4732 exemplos
totais permaneceram).

\begin{table}[h]
\centering
\caption{Fine-tuning completo: resultado final no dev set completo (289--293 exemplos, sem o filtro de tamanho de grafo aplicado ao treino).}
\label{tab:finetuning-full}
\begin{tabular}{lrrr}
\toprule
\textbf{Configuração} & \textbf{exact match} & \textbf{exec} & \textbf{eval\_loss} \\
\midrule
Zero-shot (referência) & 9,0\%  & 8,87\% & -- \\
Fine-tuning completo    & \textbf{0,0\%} & \textbf{0,34\%} & 4,2321 \\
\bottomrule
\end{tabular}
\end{table}

O resultado é \textbf{pior que o zero-shot}, apesar do \emph{loss} de
treino ter caído de forma saudável e monotônica até 0,469 ao final das 3
épocas. A causa raiz identificada: o modelo nunca viu nenhum exemplo de
\texttt{oncomx} durante o treino (excluído por VRAM), mas o dev set de
avaliação final inclui \texttt{oncomx} normalmente -- o desempenho ruim
nesse domínio nunca visto arrasta a média agregada para baixo, revertendo
o ganho que o fine-tuning trouxe em \texttt{cordis}/\texttt{sdss}. Esse é
um achado metodológico central deste trabalho: \textbf{fine-tuning
completo sob restrição severa de VRAM que exclui um dos domínios-alvo pode
piorar a generalização agregada}, mesmo quando todos os indicadores de
treino (perda) sugerem que o ajuste está funcionando.

\subsection{Fine-tuning \emph{few-shot}}
\label{sec:finetuning-fewshot}

Para tentar obter alguma representação de \texttt{oncomx} sem incorrer no
mesmo custo proibitivo de VRAM ao longo de um treino completo, foi montado
um segundo experimento com uma amostra pequena e deliberadamente balanceada
entre os três bancos: 140 exemplos de \texttt{cordis} + 140 de
\texttt{sdss} (amostragem aleatória) + 20 de \texttt{oncomx} (os 20 exemplos
de menor grafo disponíveis, 1022--1025 nós -- ver limitação (iii)), 300
exemplos no total, por 10 épocas (90 passos de otimização).

\begin{table}[h]
\centering
\caption{Fine-tuning few-shot: resultado final no dev set completo.}
\label{tab:finetuning-fewshot}
\begin{tabular}{lrrr}
\toprule
\textbf{Configuração} & \textbf{exact match} & \textbf{exec} & \textbf{eval\_loss} \\
\midrule
Zero-shot (referência)  & 9,0\%  & 8,87\% & -- \\
Fine-tuning completo     & 0,0\%  & 0,34\% & 4,2321 \\
Fine-tuning few-shot     & \textbf{0,0\%} & \textbf{0,0\%} & 4,8754 \\
\bottomrule
\end{tabular}
\end{table}

O resultado também ficou abaixo do zero-shot, e abaixo até do fine-tuning
completo no exec. Diferente do fine-tuning completo, aqui não há exclusão
sistemática de um banco (os três estão representados, ainda que
\texttt{oncomx} de forma mínima) -- a causa mais provável é volume de
dados insuficiente (300 exemplos) combinado com uma limitação metodológica
própria deste experimento, descrita em (iv) abaixo: o checkpoint final
salvo corresponde a apenas 52\% do treino planejado, não ao ponto de
convergência real.

\subsection{Correção do \emph{gradient checkpointing} e fine-tuning completo revisado}
\label{sec:gc-fix}

A limitação (iii) acima -- exclusão total de \texttt{oncomx} do treino por
falta de VRAM -- motivou uma segunda iteração deste trabalho, buscando
incluir os três domínios sem reduzir \texttt{max\_source\_length}. O
mecanismo natural para isso é o \emph{gradient checkpointing} do PyTorch
(\texttt{torch.utils.checkpoint}): em vez de manter em memória todas as
ativações intermediárias da rede para a retropropagação, ele as recalcula
sob demanda durante o \emph{backward pass}, trocando tempo de computação por
memória. A flag já existia em \texttt{configs/train.json}, mas estava sem
efeito para o modelo com RGAT -- dois problemas distintos, ambos
corrigidos:

\begin{itemize}
\item A classe \texttt{Model} em \texttt{seq2seq/models/graphix/rgat.py} é
um invólucro (\emph{wrapper}) que guarda o T5 com RGAT de fato como um
atributo interno (\texttt{self.pretrain\_model}), em vez de herdar sua
hierarquia de camadas. O método padrão
\texttt{gradient\_checkpointing\_enable()} do HuggingFace percorre
recursivamente as submódulos de uma classe para ativar o checkpointing em
cada um -- como o RGAT vive dentro de um atributo e não na hierarquia
percorrida, a ativação nunca chegava lá. Corrigido delegando
explicitamente a chamada ao \texttt{pretrain\_model} interno.
\item Mesmo com a delegação corrigida, o \texttt{T5Stack.forward()} (em
\texttt{seq2seq/models/modeling\_t5.py}) constrói, quando o checkpointing
está ativo, uma função de encapsulamento (\emph{closure}) para cada bloco
que é passada a \texttt{torch.utils.checkpoint.checkpoint()} -- e essa
closure não repassava \texttt{graph\_batch}/\texttt{relation\_emb} ao
bloco, fazendo o RGAT ser silenciosamente ignorado (ou falhar) sempre que
o checkpointing estava ativo. Corrigido capturando esses dois objetos por
\emph{closure} do Python (não são tensores diferenciáveis que o mecanismo
de checkpoint precise gerenciar, então isso é seguro tanto na chamada
inicial quanto no recálculo do \emph{backward}).
\end{itemize}

A correção foi validada com um teste de fumaça isolado sobre os seis
maiores grafos de \texttt{oncomx} (1217--1308 nós, os que mais
frequentemente causavam OOM antes), confirmando treino sem erros e perda
decrescente normalmente.

Com o \emph{gradient checkpointing} funcionando, o fine-tuning completo foi
refeito com os mesmos hiperparâmetros da Seção~\ref{sec:finetuning-full}
(3 épocas, \texttt{batch\_size=1}, \texttt{gradient\_accumulation\_steps=32},
$5\times10^{-5}$ com \texttt{adafactor}), desta vez sem nenhuma exclusão de
exemplo por tamanho de grafo -- os 4732 exemplos de treino (incluindo
\texttt{oncomx} por completo) participaram do treino.

\begin{table}[h]
\centering
\caption{Fine-tuning completo, antes e depois da correção do gradient checkpointing.}
\label{tab:finetuning-full-gc}
\begin{tabular}{lrrrl}
\toprule
\textbf{Configuração} & \textbf{exact match} & \textbf{exec} & \textbf{eval\_loss} & \textbf{train\_loss / tempo} \\
\midrule
Fine-tuning completo (sem \texttt{oncomx}) & 0,0\% & 0,34\% & 4,2321 & 0,469 / -- \\
Fine-tuning completo (com \texttt{oncomx}, gc) & 0,0\% & \textbf{0,34\%} & \textbf{1,6765} & 2,215 / 9h42min \\
\bottomrule
\end{tabular}
\end{table}

O resultado é um achado genuinamente nuançado: incluir \texttt{oncomx}
reduziu o \texttt{eval\_loss} em quase $2{,}5\times$ (4,2321 $\to$ 1,6765) --
uma melhora clara e esperada, já que o modelo agora vê exemplos do domínio
que antes lhe era totalmente estranho. Mas essa melhora \textbf{não se
traduziu em nenhum ganho nas métricas discretas}: \texttt{exact\_match}
permanece em 0\% e \texttt{exec} praticamente inalterado (0,34\%). Isso
sugere que a hipótese da Seção~\ref{sec:finetuning-full} (exclusão de
\texttt{oncomx} como causa única do resultado ruim) era \emph{necessária
mas não suficiente}: mesmo com todos os três domínios representados no
treino, os erros do modelo continuam sendo do tipo que quebra
\texttt{exact\_match}/\texttt{exec} de forma binária -- nome de
tabela/coluna incorreto, estrutura de \texttt{JOIN} equivocada -- e uma
perda agregada mais baixa (que mede o quão "próxima" a distribuição de
tokens gerada está do alvo, token a token) não impede esse tipo de erro
categórico. As causas mais prováveis são um orçamento de treino ainda
pequeno (3 épocas sobre 4732 exemplos heterogêneos, contra 5 épocas sobre
8577 exemplos homogêneos no Spider original) e/ou uma limitação de
capacidade do próprio t5-base diante da diversidade de esquemas do
ScienceBenchmark, discutida em conjunto com o achado da próxima seção.

\subsection{Baseline t5-small sem RGAT}
\label{sec:t5small-baseline}

Para isolar o efeito específico do RGAT dos demais fatores (tamanho do
modelo, orçamento de treino, dados), foi conduzido um experimento de
ablação: um T5 padrão, \textbf{sem nenhuma camada RGAT ou estrutura de
grafo}, na variante \texttt{t5-small} (60M parâmetros, contra os $\sim$220M
do \texttt{t5-base} usado em todos os experimentos anteriores). O restante
do pipeline foi mantido idêntico -- mesmo pré-processamento (a mesma
serialização de esquema textual produzida pelas etapas de grafo, ainda que
o grafo em si não seja consumido), mesmo split de treino/dev (4732/299
exemplos), mesmos hiperparâmetros (3 épocas, \texttt{batch\_size=1},
\texttt{gradient\_accumulation\_steps=32}, $5\times10^{-5}$ com
\texttt{adafactor}). Uma nova classe \texttt{Model}
(\texttt{seq2seq/models/graphix/plain.py}) carrega um
\texttt{T5ForConditionalGeneration} padrão do HuggingFace e ignora
completamente o \texttt{graph\_batch}, selecionável por uma variável de
ambiente (\texttt{GRAPHIX\_MODEL\_VARIANT=plain}) sem duplicar o restante
do script de treino.

\begin{table}[h]
\centering
\caption{t5-small sem RGAT vs.\ t5-base com RGAT (fine-tuning completo).}
\label{tab:t5small-vs-rgat}
\begin{tabular}{lrrrl}
\toprule
\textbf{Configuração} & \textbf{exact match} & \textbf{exec} & \textbf{eval\_loss} & \textbf{tempo de treino} \\
\midrule
t5-base + RGAT (com \texttt{oncomx}, gc) & 0,0\% & 0,34\% & 1,6765 & 9h42min \\
t5-small \textbf{sem} RGAT (3 épocas)     & \textbf{1,38\%} & \textbf{1,37\%} & \textbf{1,1825} & \textbf{2h53min} \\
\bottomrule
\end{tabular}
\end{table}

O resultado é surpreendente: o modelo \emph{menor} e \emph{sem} qualquer
estrutura de grafo superou o modelo com RGAT em todas as métricas -- exact
match 4$\times$ maior (1,38\% vs.\ 0,34\% de exec no RGAT, comparação
direta de exec: 1,37\% vs.\ 0,34\%, também $\sim$4$\times$), \texttt{eval\_loss}
menor, e treino $3{,}4\times$ mais rápido (sem o custo de recomputação do
\emph{gradient checkpointing} nem a sobrecarga de propagar mensagens em
grafos de até 1300 nós por exemplo). É importante notar que essa
comparação \textbf{não é perfeitamente controlada}: t5-small e t5-base+RGAT
diferem tanto na presença do RGAT quanto no tamanho do modelo base, então
o experimento isola "RGAT vs.\ sem RGAT" apenas de forma aproximada, não
com uma variável por vez. Ainda assim, o resultado é informativo na direção
oposta à esperada -- se o RGAT estivesse agregando valor líquido, seria
razoável esperar que o modelo \emph{maior e com estrutura adicional}
superasse o menor e mais simples, não o contrário. Hipóteses para essa
inversão:

\begin{itemize}
\item O RGAT deste projeto foi desenhado e ajustado para os esquemas
pequenos e homogêneos do Spider ($\sim$5,3 tabelas/banco); grafos de
1000+ nós do ScienceBenchmark estão 2--3 ordens de grandeza acima da escala
para a qual a arquitetura foi validada, e a propagação de mensagens em
grafos tão densos pode diluir sinal útil em vez de reforçá-lo.
\item Um modelo maior e mais complexo (t5-base+RGAT) tipicamente precisa de
mais dados e/ou mais épocas para superar um modelo menor em fine-tuning --
com apenas 4732 exemplos e 3 épocas, o orçamento de treino pode ser
insuficiente para que a capacidade extra do RGAT se traduza em vantagem
prática, enquanto o t5-small, por ser mais simples, converge mais
rapidamente para um padrão razoável com o mesmo orçamento.
\item A complexidade adicional (mais parâmetros, mais uma fonte de
gradiente via o \emph{encoder} do grafo) pode estar introduzindo ruído de
otimização que atrapalha a convergência dentro do orçamento de treino
disponível, sem chegar a comprometer a estabilidade geral (o \emph{loss}
de treino do RGAT+t5-base ainda caiu de forma saudável, só não tão baixo
quanto seria necessário para competir com o t5-small nas métricas
discretas).
\end{itemize}

Este resultado não permite concluir que RGAT seja inerentemente inadequado
para text-to-SQL sobre esquemas grandes -- apenas que, \textbf{nas
condições específicas deste trabalho} (hardware de 8\,GB de VRAM, poucas
horas de treino disponíveis, 4732 exemplos), a complexidade adicional do
RGAT não se pagou em métricas discretas frente a um baseline T5 puro e
menor, treinado com o mesmo orçamento de tempo.

\subsection{Limitações metodológicas}
\label{sec:limitations}

Os dois experimentos de fine-tuning expuseram uma série de limitações,
tanto de hardware quanto de configuração, documentadas aqui por
transparência e como recomendação para trabalho futuro:

\begin{enumerate}
\item[(i)] \textbf{Ausência de split de teste genuíno.} O ScienceBenchmark,
como distribuído, só define \texttt{train} (seed+synth) e \texttt{dev}. O
mesmo split de \texttt{dev} usado para validação durante o treino
(seleção do "melhor" checkpoint) foi reutilizado como avaliação final --
não há um terceiro split, cego durante todo o processo, para reportar como
número de teste independente. O Spider, usado no treino-base, tem a mesma
limitação estrutural (sem test set público), então isso é consistente com
a metodologia herdada, mas permanece uma ressalva válida: o número final
reportado pode estar levemente inflado em relação a um teste verdadeiramente
cego.

\item[(ii)] \textbf{\texttt{torch.cuda.empty\_cache()} incondicional.} Uma
chamada incondicional a \texttt{torch.cuda.empty\_cache()} em
\texttt{seq2seq/models/graphix/rgat\_tuning.py::RGAT\_Layer.forward}, uma
vez por camada RGAT (12$\times$ por exemplo no t5-base), forçava uma
sincronização CUDA completa a cada chamada. Esse código é compartilhado
entre treino e avaliação (\texttt{modeling\_t5.py} importa a mesma classe
para ambos), então os resultados de zero-shot/PICARD/beam4 relatados acima
provavelmente rodaram mais devagar do que precisavam -- os números
continuam válidos, só o tempo de execução foi maior que o necessário. O
fix (só chamar \texttt{empty\_cache()} sob OOM real, não preventivamente)
foi validado por comparação direta: antes, nenhum passo de otimização
completava em mais de uma hora; depois, o mesmo passo completava em
minutos.

\item[(iii)] \textbf{\texttt{oncomx} tem esquema uniformemente grande
demais para 8\,GB de VRAM.} Todos os 1165 exemplos de treino de
\texttt{oncomx} produzem grafos RGAT entre 1022 e 1308 nós -- não existe
um subconjunto pequeno e seguro para amostrar. Mesmo os 20 exemplos
"menores" usados no few-shot (1022--1025 nós) ainda estão acima do teto de
900--1000 nós que se mostrou confiável durante o treino completo, e
provocaram OOMs individuais recorrentes durante o few-shot (capturados e
pulados por um mecanismo de segurança, não por acaso -- ver item v). Esta é
uma restrição de hardware, não uma limitação de amostragem: nenhuma
estratégia de seleção de exemplos resolve o problema sem reduzir
\texttt{max\_source\_length} (o que exigiria reprocessar toda a base de
grafos RGAT, fora do escopo deste trabalho). \textbf{Atualização}: esta
restrição foi contornada numa iteração posterior deste trabalho habilitando
corretamente o \emph{gradient checkpointing} do PyTorch, o que permitiu
incluir \texttt{oncomx} por completo no treino sem reduzir
\texttt{max\_source\_length} -- ver Seção~\ref{sec:gc-fix}.

\item[(iv)] \textbf{Ausência de \emph{warmup} de taxa de aprendizado no
fine-tuning completo, e espaçamento de checkpoint excessivo no few-shot.}
O fine-tuning completo usou \texttt{warmup\_ratio=0}, fazendo a taxa de
aprendizado saltar diretamente para $5\times10^{-5}$ no primeiro passo --
isso corresponde a um colapso temporário de geração (repetição de
n-gramas, ex.\ \texttt{"...loadversion = 301 and loadversion = 301 and..."}
até o limite de tokens), causando \texttt{eval\_exact\_match=0} em todos os
checkpoints intermediários mesmo com o \emph{loss} caindo normalmente. Um
\texttt{no\_repeat\_ngram\_size=3} foi adicionado durante o treino
(reduzindo o tempo de geração em $5\times$ e eliminando os loops de
n-grama), e o few-shot já nasceu com \texttt{warmup\_ratio=0,1} e
\texttt{no\_repeat\_ngram\_size=3} desde o início -- mas um padrão de
repetição mais sutil (alternância entre tokens SentencePiece de tamanhos
diferentes que decodificam para o mesmo caractere repetido, ex.\
\texttt{"6666666..."}) persistiu, sugerindo que o modelo continuou
sub-treinado para geração de literais numéricos mesmo ao final do processo.
Separadamente, no few-shot, o intervalo entre checkpoints
(\texttt{eval\_steps=47}) foi ajustado no meio do treino para reduzir o
tempo gasto em avaliação (cada rodada de eval levava até 111 minutos), mas
isso fez com que só um checkpoint intermediário (passo 47 de 90, 52\% do
treino) fosse salvo antes do fim -- o \texttt{load\_best\_model\_at\_end}
carregou esse checkpoint como "melhor" por ser o único candidato com
métrica registrada, e os pesos do passo 90 (o treino de fato mais avançado)
nunca chegaram a ser persistidos em disco.

\item[(v)] \textbf{Seleção de "melhor checkpoint" quebrada quando a métrica
empata.} O fine-tuning completo usou
\texttt{metric\_for\_best\_model="exact\_match"}. Como o colapso de geração
descrito em (iv) manteve \texttt{exact\_match=0,0} em \emph{todos} os
checkpoints intermediários, e a lógica de comparação do HuggingFace
Trainer só substitui o "melhor" registrado quando a nova métrica é
\emph{estritamente} maior, o primeiro checkpoint avaliado (passo 50 de
330, a versão \emph{menos} treinada de todas) nunca foi superado e acabou
sendo o modelo efetivamente salvo ao final do treino -- os pesos do passo
300 (97\% do treino) precisaram ser recuperados manualmente do diretório de
checkpoint correspondente para a avaliação real reportada na
Tabela~\ref{tab:finetuning-full}. O few-shot corrigiu isso usando
\texttt{metric\_for\_best\_model="eval\_loss"} com
\texttt{greater\_is\_better=false} (uma métrica contínua, que decresceu de
forma monotônica ao longo do treino), mas essa correção por si só não
evitou a limitação (iv) relacionada ao espaçamento de checkpoints.

\item[(vi)] \textbf{OOMs genuínos e travamento de driver de GPU.} Dois
crashes de OOM (\emph{out of memory}) genuínos interromperam o fine-tuning
completo no meio do treino, e um deles coincidiu com um travamento
completo do driver NVIDIA (\texttt{nvidia-smi} parou de reconhecer a
placa), exigindo reinicialização completa da máquina -- um reset de driver
isolado não foi suficiente para restaurar o desempenho normal da GPU.
Mitigações aplicadas: \emph{checkpointing} a cada 50 passos (em vez de 500,
o intervalo original herdado do treino no Spider, que nunca dispararia
dentro de um treino de 330 passos) e uma captura de OOM por micro-batch em
\texttt{seq2seq/utils/trainer.py::Seq2SeqTrainer.training\_step} que pula
apenas o exemplo problemático (zerando o gradiente parcialmente acumulado
e liberando cache) em vez de derrubar o processo inteiro -- validado
diretamente durante o few-shot, que sobreviveu a dezenas de OOMs
individuais de exemplos grandes de \texttt{oncomx} sem interrupção.
\end{enumerate}

\section{Comparação final}
\label{sec:final-comparison}

\begin{table}[h]
\centering
\caption{Comparação de todas as configurações avaliadas no dev set do ScienceBenchmark.}
\label{tab:final-comparison}
\begin{tabular}{lrrl}
\toprule
\textbf{Configuração} & \textbf{exact match} & \textbf{exec} & \textbf{Observação} \\
\midrule
Zero-shot (greedy, num\_beams=1)  & 9,0\%  & 8,87\% & referência \\
Zero-shot + beam search (n=4)      & \textbf{9,34\%} & \textbf{9,56\%} & melhor resultado geral \\
Zero-shot + PICARD                 & 9,0\%  & 8,87\% & PICARD inerte (falha de integração, \S\ref{sec:picard-attempt}) \\
Fine-tuning completo, sem \texttt{oncomx} (3 épocas)    & 0,0\%  & 0,34\% & pior que zero-shot (\S\ref{sec:finetuning-full}) \\
Fine-tuning few-shot (300 ex., 10 épocas) & 0,0\% & 0,0\% & pior que zero-shot (\S\ref{sec:finetuning-fewshot}) \\
Fine-tuning completo, com \texttt{oncomx} + gc (3 épocas) & 0,0\% & 0,34\% & eval\_loss $2,5\times$ menor, métricas discretas inalteradas (\S\ref{sec:gc-fix}) \\
Baseline t5-small sem RGAT (3 épocas)     & 1,38\% & 1,37\% & supera todos os fine-tunings com RGAT (\S\ref{sec:t5small-baseline}) \\
\bottomrule
\end{tabular}
\end{table}

O achado central deste conjunto de experimentos permanece, em essência,
\textbf{negativo em relação ao fine-tuning com RGAT}: em nenhuma das quatro
tentativas de ajuste fino (completo sem \texttt{oncomx}, \emph{few-shot},
completo com \texttt{oncomx} via \emph{gradient checkpointing}) o modelo
t5-base+RGAT superou o desempenho do checkpoint zero-shot com busca em
feixe nas métricas discretas -- mesmo depois de resolver a limitação de
VRAM que impedia incluir todos os domínios de treino. A correção do
\emph{gradient checkpointing} (\S\ref{sec:gc-fix}) trouxe uma melhora real
e mensurável na perda de avaliação (2,5$\times$ menor), mas essa melhora
não se propagou às métricas discretas -- um resultado que, por si só, já é
instrutivo: perda agregada e acerto exato/execução podem divergir
significativamente quando os erros do modelo são predominantemente
categóricos (nome de tabela/coluna errado) em vez de graduais.

O achado mais forte deste trabalho, no entanto, veio do baseline de
ablação (\S\ref{sec:t5small-baseline}): um \texttt{t5-small} \textbf{sem
nenhuma estrutura de grafo}, treinado com o mesmo orçamento de dados e
épocas, superou \emph{todas} as configurações com RGAT fine-tunadas neste
trabalho em exact match e exec, e o fez em menos de um terço do tempo de
treino. Combinado com o resultado do zero-shot (onde o RGAT+t5-base é a
única opção avaliada, pois não há um t5-small equivalente pré-treinado no
Spider para comparação direta), a evidência deste trabalho aponta para uma
conclusão qualificada: \textbf{nas condições de hardware e dados
disponíveis, a complexidade adicional do RGAT não demonstrou vantagem
sobre um T5 padrão no fine-tuning para o ScienceBenchmark} -- um resultado
negativo genuíno para a hipótese de que a estrutura de grafo ajudaria a
generalizar para esquemas maiores e de domínio distinto do treino
original. As limitações documentadas na Seção~\ref{sec:limitations}
(volume de dados, orçamento de épocas, ausência de um terceiro split de
teste cego) continuam válidas e sugerem que essa conclusão é específica
deste orçamento experimental, não necessariamente uma propriedade
fundamental da arquitetura RGAT -- mas dentro do escopo efetivamente
testado, o T5 puro foi tanto mais eficaz quanto mais barato de treinar.
